\documentclass[final]{agujournal2019}
\usepackage{amsmath}
\linenumbers

% -------------------------------------------------------------------------
\usepackage{float}
\graphicspath{{figures/}}
\usepackage{booktabs}
\usepackage{array}
\usepackage{tabularx}
\newcolumntype{L}{>{\raggedright\arraybackslash}X}
\newcolumntype{R}{>{\raggedleft\arraybackslash}X}
\usepackage[final]{listings}
\lstset{basicstyle=\ttfamily\footnotesize,breaklines=true,frame=single,
        captionpos=b,showstringspaces=false}

\usepackage{algorithm}
\usepackage{algpseudocode}


\draftfalse

\journalname{JGR: Machine Learning and Computation}

\begin{document}

% -------------------------------------------------------------------------
\title{Resolving window dependence in flexural parameter recovery at twenty
Pacific trenches by physics-informed neural inversion}


\authors{ %
  Polina Lemenkova\affil{1},
  Alexey L. Piskarev\affil{1,2}
}

\affiliation{1}{Saint-Petersburg State University, St. Petersburg 199034, Russia}
\affiliation{2}{Institute for Geology and Mineral Resources of the Ocean
(VNIIOkeangeologia), St. Petersburg 190121, Russia}

% -------------------------------------------------------------------------
\begin{keypoints}
\item Moving only the interval fitted on a fixed bending profile changes the
recovered elastic thickness by about forty per cent
\item An elastic plate equation imposed at every point along the profile
recovers the thickness without any fitted interval
\item Recovered thickness is near uniform at nine kilometres and one fifth of
the mechanical thickness with a physics-informed neural network (PINN) across twenty Pacific trenches
\end{keypoints}

\begin{abstract}
The oceanic lithosphere bends downward before it enters a subduction zone, and
the resistance of the plate to that bending is summarised by its effective
elastic thickness, the parameter through which the mechanical state of the
incoming plate is read. Published determinations disagree, and part of the
disagreement has never been examined: every profile-based inversion is fitted
over an interval of the profile chosen by the analyst, and that interval is
imposed before the fit and reported with no uncertainty of its own. Here the
flexural parameters of twenty circum-Pacific trenches are recovered without such
an interval, by a physics-informed neural network (PINN). The deflection is
represented by a neural network, the thin-elastic-plate equation is imposed as a
residual at collocation points spanning the whole profile, and the elastic
thickness is a trainable coefficient recovered jointly with the network across 71
along-strike segments, from open-access bathymetry, seafloor age, sediment
thickness and satellite gravity. Because the governing equation enters the
training objective itself, the PINN estimate is constrained by the physics
everywhere on the profile rather than by an interval imposed upon it.
Sweeping the interval through the classical inversion, on fixed profiles and a
fixed physical model, reproduces an interquartile spread in recovered thickness
of 40 per cent of the profile median, and a full range across defensible
windows of 121 per cent, from 428 profiles at twenty margins. A measurable
share of the published disagreement is therefore procedural rather than
physical, and the reported metric allows any profile-fitting study to quantify
the same dependence on its own data.
\end{abstract}

\section*{Plain Language Summary}
Where one of the great plates of the sea floor slides beneath another, it first
bends downward, and the sea floor in front of the deep trench rises into a long,
low swell. The shape of that bend records how stiff the plate is, and stiffness
matters because it controls where the plate cracks, how those cracks let sea
water into the rocks below, and how the plate carries the forces that drive it
down. Stiffness is measured by fitting a simple physical model of a bending
sheet to the observed shape of the sea floor. Doing so requires a decision about
which stretch of sea floor to include in the fit, and there is no rule for
making that decision. Different choices give different answers from the
same sea floor. This study measures how large that effect is, and then removes
the decision altogether by requiring the bending equation to hold everywhere
along the profile instead of fitting a chosen stretch of it. Applying one
automatic procedure to twenty trenches around the Pacific Ocean, using only
freely available data, gives a set of stiffness values that can be compared with
each other for the first time. The procedure represents the bent sea floor by a
neural network that is required to satisfy the bending equation itself, an
approach known as a physics-informed neural network.

% =========================================================================
\section{Introduction}\label{sec:intro}

Subduction begins with bending. Before the oceanic lithosphere descends beneath
an overriding plate it flexes downward under the load applied at the plate
boundary, and that flexure produces the two morphological features that define
every convergent margin: the trench itself and the broad topographic swell of
the outer rise seaward of it \cite{Kirby2022Isostasy,Kirby2022Flexure}. The
resistance of the plate to this deformation is summarised by a single quantity,
the effective elastic thickness, which sets the flexural rigidity and therefore
controls the wavelength of the deflection, the amplitude and position of the
outer-rise crest, and the magnitude of the bending stresses carried by the
plate \cite{ContrerasReyesOsses2010,Manriquez2014}. Because those stresses
govern where the plate fails in extension, the effective elastic thickness also
conditions outer-rise normal faulting and, through the fractures it opens, the
pathways by which seawater penetrates the crust and uppermost mantle
\cite{ContrerasReyes2007,Craig2014}. It is, in short, the parameter through
which the mechanical state of the incoming plate is read.

Estimating that parameter has therefore been a persistent objective of marine
geophysics, and the established approach is to fit the deflection predicted by a
thin elastic plate to the observed shape of the margin. The plate is treated
either as continuous or as broken at the trench axis, and the free parameters
are adjusted until the predicted deflection reproduces the bathymetric profile,
the free-air gravity anomaly, or both together
\cite{Tandon2000,ContrerasReyesOsses2010}. Spectral formulations pursue the same
end through the admittance between topography and gravity rather than through
the profile shape \cite{CalmantCazenave1987,Manea2005}, and three-dimensional
treatments relax the assumption of a laterally uniform plate by allowing the
elastic thickness to vary along and across strike \cite{Manriquez2014}. Each of
these formulations has been made possible by the global coverage of
satellite-derived bathymetry and marine gravity, which supply profiles at
margins where shipborne surveys remain sparse
\cite{Sandwell2014,CalmantBaudry1996,BaudryCalmant1996}. The physical model is
thus shared across the literature; what differs is how it is fitted.

That difference has consequences, because the values reported for individual
margins do not agree. Estimates vary substantially from one subduction zone to
another, which is expected given the range of plate ages and loading conditions
involved \cite{DiGiuseppe2009,SchettinoTassi2011}, but they also vary between
studies of the same margin, which is not. Several physical explanations have
been advanced for the discrepancy. Inelastic yielding reduces the recovered
thickness below the mechanical thickness of the plate, so that the fitted value
depends on how much of the bending the model is asked to reproduce
\cite{ContrerasReyesOsses2010,Tandon2000}. Extensional failure at the outer rise
weakens the plate progressively as it approaches the axis, making a single
uniform value an approximation to a quantity that changes along the profile
\cite{Craig2014,KaoChen1996}. Hydration and thermal structure have both been
invoked, as has genuine along-strike variability within a single trench
\cite{ContrerasReyes2007,MoscosoContrerasReyes2012}. These explanations are
plausible and probably all contribute, yet they share an assumption worth
examining: that the disagreement is physical in origin.

A procedural source of disagreement has received far less attention. Every
profile-based inversion requires the analyst to choose an interval of the
profile over which the elastic model is fitted --- a seaward limit beyond which
the data are considered unaffected by bending, and a landward limit inside which
the elastic description is considered to fail. This analysis window is not
estimated from the data; it is imposed before the fit, and it differs between
studies because no procedure fixes it. Published work adopts windows of
different length, different origin relative to the trench axis, and different
treatment of the outer-rise crest \cite{Manriquez2014,Tandon2000}, and the
recovered thickness is reported without an uncertainty attributable to that
choice \cite{ContrerasReyesOsses2010}. The window is therefore a free parameter
of the method that enters every published value and is quantified in none of
them. Its contribution to the scatter in the literature has never been isolated
from the physical contributions described above.

Isolating it matters for interpretation rather than for tidiness. Effective
elastic thickness is not an end in itself; it is used to infer the depth of the
controlling isotherm, to estimate the degree of inelastic reduction relative to
the mechanical thickness, and to compare the strength of plates of differing
age across the ocean basins \cite{DiGiuseppe2009,CalmantCazenave1987}. Each of
those inferences propagates whatever error the fitting procedure introduces. If
a substantial share of the observed spread is procedural, then age-strength
relations and isotherm depths calibrated against the published compilation
inherit a bias whose size is unknown. The problem is sharpest for basin-wide
comparison, where values obtained under different conventions are placed on a
single plot: any such comparison presumes a common procedure that does not
exist. Applying one automated, script-driven treatment to every margin of the
Pacific basin removes that presumption, in the manner by which basin-wide
structural syntheses are now assembled from open grids
\cite{Wessel2019,Bird2003,Piskarev2019Book,Lemenkova202213,Muller2017,Sandwell2014}.

The formulation adopted here removes the window rather than standardising it.
Instead of minimising a misfit over a chosen interval, the deflection is
represented by a neural network whose derivatives are obtained through automatic
differentiation, and the governing flexure equation is imposed as a residual
evaluated at collocation points distributed across the entire profile
\cite{Raissi2019,Karniadakis2021}. The solution is constrained by the physics at
every point rather than by the extent of a fitting interval, so no seaward or
landward limit enters the estimate. Such a formulation is termed a physics-informed neural network
(PINN) \cite{Raissi2019,Karniadakis2021}. The term denotes a training
framework rather than a learning algorithm in the sense that gradient-boosted
trees or random forests are algorithms: the network architecture is
conventional, and what distinguishes the approach is the objective it is trained
against. Alongside any misfit to observations, the loss carries the residual of
the governing differential equation, evaluated by automatic differentiation at
points where no measurement exists, so that a candidate solution is penalised
for violating the physics as well as for departing from the data. Two
consequences follow, and both bear directly on the problem posed here. Where
observations are sparse or unevenly distributed the residual continues to
constrain the solution, because it can be evaluated anywhere in the domain;
and where a coefficient of the governing equation is unknown, it may be declared
trainable and recovered jointly with the network, which converts a forward
solver into an inverse one without a separate optimisation loop.

That second property is why the framework suits the present problem. The
governing equation of plate flexure is known and its coefficient, the flexural
rigidity, is precisely the quantity sought, so the physics can be imposed while
the unknown parameter is learned. The approach has developed rapidly since its
introduction and now includes gradient-enhanced and operator-learning variants
\cite{Wang2025gIBPINN,Yang2026DeepONet}, together with treatments of the
spectral bias that limits how finely such a network resolves a profile
\cite{WangWangPerdikaris2021}. Within the solid Earth sciences PINNs have been
applied most extensively to the seismic wave equation, both as forward solvers
and within full-waveform inversion
\cite{RashtBehesht2022,SongAlkhalifah2021,ZhangZhuGao2023,Schuster2024}, to the
eikonal equation for traveltime computation \cite{Waheed2021PINNeik}, to
hypocentre inversion under a differentiable traveltime model
\cite{Smith2022HypoSVI}, and to petrophysical inversion
\cite{Li2024Probabilistic,Li2024Bayesian}, in formulations that impose the
residual through fully connected networks or through convolutional and recurrent
architectures encoding the discretisation directly
\cite{Gao2021PhyGeoNet,Ren2022PhyCRNet,Ren2024SeismicNet,Ding2023Topography,Ding2023HalfSpace,Hughes2019}.
What these applications share is the setting in which a PINN is most
advantageous: the governing equation is established, the observations are
incomplete, and the quantity of interest is a coefficient or a field that
classical inversion recovers only under auxiliary choices.
A parallel line of work infers the governing equation itself from data rather
than imposing it \cite{BongardLipson2007,SchmidtLipson2009,Brunton2016,Rudy2017,CampsValls2023,YuWang2024},
and the two meet where a known equation is imposed while its coefficients are
estimated \cite{Chen2021Scarce,Rao2023}, which is the position taken here. Its application to
lithospheric flexure, where the governing equation is a single ordinary
differential equation with a physically meaningful coefficient, has not been
reported. The magnitude of the effect that motivates this reformulation is
illustrated for a single profile in Figure~\ref{fig:windowproblem}.

\begin{figure}\centering
  	\noindent\includegraphics[width=0.90\textwidth]{fig_windowproblem}
  \caption{Dependence of the recovered effective elastic thickness on the
  analysis window, for a single trench-normal profile. The two panels plot
  different quantities. \textbf{Main panel:} sampled depth against distance
  from the trench axis, positive seaward. Grey, the observed section;
  coloured, the broken-plate solutions obtained by least-squares inversion of
  the seaward slope within four alternative windows, each labelled with the
  effective elastic thickness it returns. The broken vertical line marks the
  axis; the landward reach is greyed and enters no fit. \textbf{Inset:}
  recovered effective elastic thickness against the seaward window limit,
  from one hundred and one independent inversions of the same profile at
  increments of two kilometres, with the four windows of the main panel
  marked in the corresponding colours. Software used for plotting figure:
  Python \texttt{3.12.13}, Matplotlib \texttt{3.11.1}, SciPy
  \texttt{1.18.0}, NumPy \texttt{2.4.6}. Data: GEBCO. Source: authors.}
  \label{fig:windowproblem}
\end{figure}

Four windows drawn from the published range return four different values of
the effective elastic thickness from that single profile, and sweeping the
seaward limit continuously shows the estimate drifting rather than settling on
one. Neither the observations nor the physical model changes across those
inversions; only the interval over which the model is fitted.

This study therefore pursues two coupled objectives: to measure how much of the
disagreement in published flexural parameters is attributable to the analysis
window, and to recover those parameters without one. Both are carried out on the
same set of twenty circum-Pacific trenches, using open-access grids throughout,
so that the methodological result and the geophysical result rest on identical
data. The classical windowed inversion is reproduced faithfully and swept across
a systematic grid of window limits, and the physics-informed formulation is
trained jointly across all margins with the flexural parameters as trainable
coefficients. The contributions are as follows:
\begin{itemize}
  \item a quantification of the window-induced spread in recovered flexural
  parameters, measured across twenty subduction margins under a single
  reproducible procedure;
  \item a differentiable elastic-plate formulation that recovers the effective
  elastic thickness without an analysis window, validated against synthetic
  profiles of known parameters;
  \item joint recovery of the full flexural parameter set --- elastic
  thickness, bending moment, shear force and maximum curvature --- with
  ensemble-based uncertainty propagated to each;
  \item an open implementation and a released per-segment parameter database,
  enabling every figure to be regenerated from the cited open-access grids.
\end{itemize}

The remainder of the paper is organised as follows. Section~\ref{sec:data}
describes the twenty margins and the open-access datasets from which the bending
profiles are constructed. Section~\ref{sec:methods} sets out the preprocessing
chain, the governing equation shared by both inversions, the windowed baseline,
the PINN formulation, and the validation and uncertainty design.
Section~\ref{sec:results} reports the window-sweep experiment and the recovered
parameters, Section~\ref{sec:discussion} interprets them against the published
compilation, and Section~\ref{sec:conclusions} states the conclusions.

% =========================================================================
\section{Study Area and Data}\label{sec:data}

\subsection{Study Area}\label{sec:studyarea}

The circum-Pacific belt is the only setting in which one inversion procedure can
be applied to a nearly continuous population of subducting plates spanning most
of the range of lithospheric ages and convergence rates found on Earth. Twenty
trenches are treated here (Figure~\ref{fig:studyarea}). Together they enclose
the basin from the Aleutian arc in the north, southward along the western margin
through the intra-oceanic arcs of the Philippine Sea region, across the
southwest Pacific to New Zealand, and back along the eastern margin from
Cascadia to southern Chile. Each of these margins presents the same pair of
morphological features on its seaward side: an axis of maximum depth and, beyond
it, a broad flexural swell \cite{Stern2002}. That pair is the observable which
the elastic-plate model is written to reproduce, and its presence at every one
of the twenty margins is what makes a single, uniformly applied procedure
defensible.

One of the twenty requires a qualification that is carried through the whole
analysis. At Cascadia the subduction zone is unquestionably active -- the Juan
de Fuca plate continues to descend beneath North America -- but the trench
itself has no bathymetric expression, having been buried beneath the turbidite
fill of the Astoria and Nitinat fans and the accretionary prism built from it.
The margin is therefore a sediment-buried, or morphologically extinct, trench
rather than a relict subduction zone: what has been lost is the landform, not
the process. The distinction matters here because the observable this study
inverts is the landform. Where the axial fill is thicker than the flexural
deflection it conceals, no axis of maximum depth can be traced and no outer-rise
crest can be located, so the pair of features named above is absent and the
elastic-plate model has nothing to be fitted to. Cascadia is the margin of this
belt at which that condition holds, and it is therefore absent from the
inventory of Table~\ref{tab:trenches} rather than carried through it as a row of
zeros: the belt is described here by the margins whose landform records the
bending, and the one active margin whose landform does not is named and
accounted for in this paragraph instead.

The twenty margins fall into three groups that differ in the properties
governing flexure. The northwestern and western Pacific group consumes the
oldest oceanic lithosphere on the planet beneath intra-oceanic or thin
continental upper plates, and contains the deepest trenches and the most starved
axial fills; it runs from the Aleutian arc through Kuril-Kamchatka, Japan,
Izu-Bonin, Mariana, Yap and Palau to Ryukyu, Manila and the Philippine Trench.
The southwestern group, from New Britain and San Cristobal through Vanuatu,
Tonga and Kermadec to Hikurangi and Puysegur, is characterised by rapid
convergence, active trench rollback and young back-arc basins, so that the upper
plate exerts a stronger control on trench position than elsewhere. The eastern
group is of Andean type, where the young and warm lithosphere of the Juan de
Fuca, Cocos and Nazca plates descends beneath continental margins at Cascadia,
Middle America and Peru-Chile, and where the terrigenous supply is at some
latitudes large enough to bury the flexural signal entirely -- completely so at
Cascadia, which for that reason supplies no profile to this study.
Plate age, convergence rate and back-arc stress regime vary systematically
between the three groups \cite{HeuretLallemand2005,DeMets2010,Bird2003}, so the
set covers, rather than samples, the parameter range over which flexural
rigidity is expected to differ. The same basin has supplied the reference
determinations of trench flexure against which any new estimate is read
\cite{LevittSandwell1995,Bodine1981}, and the spectral and spatial estimators
applied elsewhere to continental lithosphere
\cite{StewartWatts1997,WieczorekSimons2005,Zamani2014} return values that are
not directly comparable with profile fits, which is part of why a single
uniform procedure is needed.

An entire trench is too coarse a unit for this purpose, because the age of the
incoming plate, the state of the outer rise and the thickness of the axial fill
all change along strike, and a parameter fitted to a whole margin averages over
those changes. The twenty traced axes total 32,364~km, and each is therefore
partitioned into bins of uniform along-strike length, 148~km, giving 220
segments; every flexural parameter reported below is a per-segment quantity.
Segmentation serves two ends at once. It confines each estimate to a reach over
which the plate can be treated as laterally uniform, and it raises the number of
independent determinations from twenty to 220, so that the window-induced spread
of Section~\ref{sec:res:window} is measured over a population large enough to
separate from geological variation. The boundaries follow the automatically
traced axis at fixed intervals rather than published province divisions, so no
interpretive segmentation of the margins enters the analysis. Axis lengths,
segment counts, retained profile numbers and mean plate ages are given per
trench in Table~\ref{tab:trenches}.

Three conventions govern that inventory. Lengths are measured along the traced
axis and not along the plate boundary, which is the longer of the two wherever
they diverge. The segment count is the partition the axis length implies at the
fixed 148~km spacing, and is therefore an upper bound on the number of segments
that carry a flexural estimate, a segment left empty after profile screening
being dropped rather than filled; the number retained is reported with
the results. Mean plate age is the arithmetic mean over the retained profiles of
each margin, sampled at the axis node, and the profile counts follow the
screening criteria of Section~\ref{sec:preprocessing}. Cascadia is not listed.
The turbidite fill of the Astoria and Nitinat fans leaves no depth minimum to
trace and no outer rise to fit, so the margin supplies no profile, carries no
segment and enters no result; a census row of zeros would state only that the
method cannot be applied there, which is said here instead. Vityaz is listed
in its place. It is a relict trench: subduction has ceased, so it has no active
plate boundary to seed the trace and its axis is seeded from the bathymetry
instead, by the same procedure that then snaps every node to the local depth
minimum. Including it keeps the inventory to the margins that contribute a
profile, and adds the one end-member in the belt where the outer rise formed
under a convergence that has since stopped.

\begin{figure}\centering
  	\includegraphics[width=0.95\textwidth]{fig_studyarea}
  \caption{The twenty circum-Pacific trenches treated in this study. Shaded
  bathymetric relief on an equidistant cylindrical projection centred on the
  date line, illuminated from the northwest. Red, the trench axes traced
  automatically as the locus of deepest nodes; black, the plate boundaries of
  the PB2002 model; thin grey, contours of seafloor age at 20 million year
  intervals; tick marks along each axis, the boundaries of the 148 kilometre
  along-strike segments. Trenches are labelled by name and grouped by the
  tectonic setting of Table~\ref{tab:trenches}. Grey, land from the shoreline
  database. Software used for plotting figure: GMT \texttt{6.5.0}. Data: GEBCO bathymetry, EarthByte seafloor age, PB2002 plate
  boundaries. Source: authors.}
  \label{fig:studyarea}
\end{figure}

\subsection{Data and Materials}\label{sec:datamaterials}

Every input is an open-access global grid or vector database; no fieldwork,
proprietary survey or restricted product enters the study, and each margin is
treated with the same release of the same source, so that no part of the
between-trench variation can be attributed to a change of data. The primary
observable is bathymetry, taken from the General Bathymetric Chart of the Oceans
(GEBCO) grid at 15 arc-second spacing \cite{GEBCO2023}. Between 50 degrees south
and 60 degrees north that grid rests on version 2.5.5 of the SRTM15+ compilation
\cite{Tozer2019}, augmented with the regional multibeam data sets of the Seabed
2030 project. Where soundings are absent the base predicts depth from satellite
gravity through a gravity-to-topography transfer \cite{Tozer2019,Sandwell2014},
so part of the deep-ocean grid is an interpolation rather than a measurement.
Later releases replace that transfer with a trained neural network
\cite{HarperSandwell2024,GEBCO2024,GEBCO2025}, which changes the character of
the predicted nodes but not the fact that they are predicted. The accompanying type identifier grid, distributed at
the same spacing, records which of the two each node is, and it is retained
alongside the depths so that profiles resting largely on predicted values can be
flagged.

Three further grids supply the quantities needed to convert a raw depth section
into a deflection and to interpret the recovered parameters. Oceanic crustal age
is taken from the EarthByte present-day compilation of Seton et al. at
2 arc-minute spacing \cite{Seton2020}, on the GTS2012 timescale
\cite{Seton2020}; it fixes the thermal reference depth against which the bending
of each profile is measured, and provides the age axis of the analysis in
Section~\ref{sec:disc:physics} \cite{Lemenkova202112}. Total sediment thickness
is taken from GlobSed
at 5 arc-minute spacing \cite{Straume2019} and supports the correction that
removes the axial fill from the observed depth, without which a well-fed trench
would be read as less deeply bent than it is. Each grid is checked for extent,
registration and units before any of it is read, so that a mismatch of
convention is caught at the source rather than propagating silently into the
deflections (Listing~\ref{lst:grids}). \begin{lstlisting}[caption={Verifying extent, registration and units of the
three source grids before any profile is extracted. The age grid is stored on
$-180/180$ while the traced axes are written on $0/360$; the discrepancy is
reconciled at sampling time rather than assumed away.},label={lst:grids},language=bash]
gmt grdinfo GEBCO_2023.nc -C | cut -f2-7   # xmin xmax ymin ymax zmin zmax
gmt grdinfo age.2020.1.GTS2012.2m.nc       # z in Ma, gridline-registered
gmt grdinfo GlobSed-v3.nc -C | cut -f2-5   # sediment thickness, m
\end{lstlisting}

Free-air gravity anomalies are
taken from the satellite-altimetry compilation of the Scripps Institution of
Oceanography at 1 arc-minute spacing \cite{Sandwell2014,CalmantBaudry1996} and
serve as an independent expression of the same deflection at margins where the
bathymetric grid is dominated by predicted values, following the long-standing
practice of reading submarine basin structure from bathymetry and satellite
gravity together
\cite{Sandwell2014,Tozer2019,Piskarev2019Makarov,BaudryCalmant1996,LevittSandwell1995}. The coverage of the three
over the study domain is shown in Figure~\ref{fig:datacoverage}.

Three vector databases complete the input set. Plate boundary geometry follows
the PB2002 model \cite{Bird2003} and assigns each segment a subducting and an
overriding plate. The MORVEL angular velocities \cite{DeMets2010} then give the
convergence rate and its obliquity at each segment; both inform the
interpretation and neither enters the inversion. Coastlines for the map are
drawn from the Global Self-consistent, Hierarchical, High-resolution Geography
(GSHHG) database \cite{WesselSmith1996}, which is also the shoreline source used
internally by the mapping software \cite{Wessel2019}. Table~\ref{tab:data} lists
each input with its variable, native spacing, release identifier and originating
publication. All are distributed without access restriction, and all were
obtained in the releases stated there; persistent identifiers are given in the
Open Research Section.

The inputs differ in spacing by more than an order of magnitude, and the
analysis is designed around that difference. At 15 arc
seconds the bathymetric grid samples the interval between the axis and the
outer-rise crest, typically fifty to one hundred and fifty kilometres, with one
hundred or more nodes, so the deflection is resolved far more finely than the
wavelength being fitted. A single cell of the sediment and age compilations, by
contrast, spans roughly a tenth of that same interval. Those two grids are
therefore never interpolated to profile resolution and applied node by node.
Each is sampled once per profile, at the axis node, and the resulting value is
held constant along that profile; age and sediment thickness are treated as
attributes of the margin at that point rather than of the individual depth
samples. The restriction keeps the sediment and thermal corrections from
imposing structure at wavelengths their source grids cannot support, which would
otherwise be absorbed into the recovered rigidity and misread as a property of
the plate.

\begin{table}\centering
  \caption{Input data sets, with the variable each supplies, its native grid
  spacing or vector form, the release used and the publication describing it.
  All are open access and were obtained in the releases stated; no restricted or
  unpublished data enter the study. Grid spacings are quoted in arc seconds and
  arc minutes as distributed. Source: authors.}
  \label{tab:data}
  \small
  \begin{tabular}{@{}llll>{\raggedright\arraybackslash}p{0.23\textwidth}@{}}
    \toprule
    Data set & Variable & Spacing & Release & Source \\
    \midrule
    GEBCO grid & Bathymetry & 15 arc-sec & GEBCO\_2023 & \citeA{GEBCO2023} \\
    GEBCO type grid & Node type & 15 arc-sec & GEBCO\_2023 & \citeA{GEBCO2023} \\
    SRTM15+ & Base grid of the above & 15 arc-sec & v2.5.5 & \citeA{Tozer2019} \\
    EarthByte age grid & Crustal age & 2 arc-min & 2020.1 & \citeA{Seton2020} \\
    GlobSed & Sediment thickness & 5 arc-min & 2019 & \citeA{Straume2019} \\
    Marine gravity & Free-air anomaly & 1 arc-min & v33.1 & \citeA{Sandwell2014} \\
    PB2002 & Plate boundaries & Vector & 2003 & \citeA{Bird2003} \\
    MORVEL & Angular velocities & Vector & 2010 & \citeA{DeMets2010} \\
    GSHHG & Shoreline polygons & Vector & v2.3.7 & \citeA{WesselSmith1996} \\
    \bottomrule
  \end{tabular}
\end{table}

\begin{figure}\centering
  	\includegraphics[width=0.95\textwidth]{fig_datacoverage}
  \caption{Coverage of the auxiliary grids over the circum-Pacific
  domain, with the traced trench axes drawn in black on every panel.
  (a) Seafloor age; (b) sediment thickness; (c) free-air gravity anomaly.
  Grey, land from the GSHHG shoreline database; white, nodes at which the
  grid carries no value. All three are shown on an equidistant cylindrical
  projection over the same region, so that coverage can be compared panel by
  panel and any margin lacking an auxiliary value can be identified before
  the parameterisation is run. The colour scales are perceptually uniform and
  colour-blind safe, and the diverging scale of panel (c) is anchored on zero.
  Software used for plotting figure: GMT \texttt{6.5.0}. Data: EarthByte
  seafloor age, GlobSed sediment thickness, IGPP free-air anomaly. Source:
  authors.}
  \label{fig:datacoverage}
\end{figure}

\begin{table}\centering
  \caption{The twenty circum-Pacific trenches, grouped by tectonic setting, with
  the length of the automatically traced axis, the number of 148~km along-strike
  segments into which it is partitioned, the number of bending profiles retained
  after screening, and the mean age of the subducting plate at the axis. The
  conventions governing each column are stated in Section~\ref{sec:studyarea}.
  Source: authors.}\label{tab:trenches}
  \begin{tabularx}{\textwidth}{@{}>{\hsize=1.60\hsize}L
                                *{4}{>{\hsize=0.85\hsize}R}@{}}
    \toprule
    Trench & Length (km) & Segments & Profiles & Age (Ma) \\
    \midrule
    \multicolumn{5}{@{}l}{\textit{Northwestern and western Pacific}} \\
    Aleutian         &  3772 &  25 & 187 &  55.7 \\
    Kuril-Kamchatka  &  2223 &  15 & 125 & 113.6 \\
    Japan            &   793 &   5 &  79 & 134.3 \\
    Izu-Bonin        &  1128 &   8 &  99 & 137.9 \\
    Mariana          &  1867 &  13 & 179 & 150.2 \\
    Yap              &   691 &   5 &  73 & 136.8 \\
    Palau            &   848 &   6 &  14 &  33.7 \\
    Ryukyu           &  1125 &   8 &  66 &  69.5 \\
    Manila           &   610 &   4 &  56 &  29.5 \\
    Philippine       &  1634 &  11 & 150 &  72.7 \\
    \addlinespace
    \multicolumn{5}{@{}l}{\textit{Southwestern Pacific}} \\
    New Britain      &   655 &   4 &  66 &  39.9 \\
    San Cristobal    &  1386 &   9 &  59 &  38.7 \\
    Vanuatu          &  1351 &   9 & 117 &  37.0 \\
    Tonga            &  1126 &   8 & 111 & 100.7 \\
    Kermadec         &  1590 &  11 & 149 & 111.1 \\
    Hikurangi        &   716 &   5 &   6 & 128.2 \\
    Puysegur         &   834 &   6 &  22 &  25.9 \\
    Vityaz\textsuperscript{$\dagger$} & 1415 & 10 & 3\textsuperscript{$\ddagger$} & --- \\
    \addlinespace
    \multicolumn{5}{@{}l}{\textit{Eastern Pacific}} \\
    Middle America   &  2256 &  15 & 163 &  15.5 \\
    Peru-Chile       &  6344 &  43 & 143 &  44.5 \\
    \midrule
    \textit{Total}   & 32{,}364 & 220 & 1867 & \\
    \bottomrule
  \end{tabularx}

  \vspace{2pt}
  {\footnotesize\raggedright
    $\dagger$ Relict trench: subduction ceased, no active plate boundary, so
    the axis is seeded from the bathymetry rather than from PB2002. No mean
    plate age is quoted, the age compilation returning no value at its
    retained profiles.\par
    $\ddagger$ Three profiles survive screening, the smallest sample of any
    margin in the inventory.\par}
\end{table}

% =========================================================================
\section{Methodology}\label{sec:methods}

\subsection{Data Preprocessing}\label{sec:preprocessing}

The trench axis is traced automatically, because a hand-digitised axis would
reintroduce at the first step the analyst dependence the study sets out to
remove. The subduction segments of the PB2002 model \cite{Bird2003} serve only
as a seed that locates the margin: they are clipped to a bounding box for each
trench, pieces shorter than forty kilometres are discarded, the remainder are
chained by nearest endpoint without bridging any gap wider than two hundred and
fifty kilometres, and the chained seed is resampled at ten-kilometre intervals.
At every seed node a search profile of half-length sixty kilometres is cast
perpendicular to the local strike and the bathymetric grid is sampled along it
at half-kilometre steps; the deepest sample becomes the axis node. The axis
therefore follows the bathymetry, and the published plate boundary constrains
only where the search begins. Three margins carry no PB2002 subduction segment
and are seeded from a short hand-digitised deep instead, after which the
snapping step proceeds identically.

Two guards protect the trace. A trench shoals gradually towards its ends, so a
global depth cut-off deep enough to exclude a parallel back-arc depression would
amputate those terminations; the discriminator used instead is the shape of the
along-strike depth series, and a node lying more than two thousand metres
shallower than the rolling median of its twenty-one neighbours is rejected as an
isolated excursion. This is what separates the Ryukyu Trench from the Okinawa
Trough, which runs parallel to it and lies inside any bounding box that holds
the trench. The surviving nodes are median-filtered over twenty-five kilometres
so that the local strike, and hence the profile azimuth, is stable. Each axis is
then measured for path length against end-to-end distance, and a ratio above 1.6
rejects the trace outright: arcuate margins reach 1.3 legitimately, so a higher
ratio means the path doubles back and would extract profiles twice along the
same margin from opposite limbs. The axis so obtained is a property of the
bathymetric grid, not of the operator, and rerunning the routine reproduces it
exactly (Listing~\ref{lst:axes}). The traced file is inspected before use, both
for the number of axis records and for the path-to-chord ratio that rejects a
doubled-back trace (Listing~\ref{lst:axischeck}).

\begin{lstlisting}[caption={Inspecting the traced axes: one record per trench,
and the along-path node count that fixes the number of profiles each margin can
carry.},label={lst:axischeck},language=bash]
grep -c '^>' trench_axes_full.gmt          # 20 axis records
grep '^>' trench_axes_full.gmt | head -3   # > -L"Aleutian"
# nodes per axis, i.e. 5 km steps along strike
awk '/^>/{if(n)print t, n; t=$0; n=0; next}{n++}END{print t, n}' \
    trench_axes_full.gmt
\end{lstlisting}

Bending profiles are extracted normal to the local tangent of the traced axis at
ten-kilometre intervals along strike, each of half-length two hundred and fifty
kilometres on either side of the axis, and each sampled at one-kilometre steps
by bilinear interpolation of the bathymetric grid. Distances and azimuths are
geodesic on the WGS84 ellipsoid, so a profile is perpendicular to the true local
strike rather than to its image in a projected plane. The seaward
half-length is set by the requirement that the profile reach essentially
undeflected seafloor for the stiffest plates expected, since a profile
terminating inside the flexural signal would impose a truncation of exactly the
kind under examination; it also fixes the outer bound of the window grid swept
in Section~\ref{sec:experiments}. The landward half is retained for display and
for locating the axis, but enters no inversion.
Figure~\ref{fig:profilegeometry} defines the geometry and the symbols used
throughout.

\begin{figure}\centering
  \includegraphics[width=0.90\textwidth]{fig_profilegeometry}
  \caption{Geometry of a bending profile and the symbols used throughout.
  A profile is cast normal to the local tangent of the traced axis and
  sampled at one-kilometre steps; distance $x$ is measured from the trench
  axis and is positive seaward, and the deflection $w$ of
  Equation~\eqref{eq:deflection} is positive downward. Marked on the section
  are the axis itself, the outer-rise crest at distance $x_{b}$ carrying
  amplitude $w_{b}$ above the reference depth, the reference depth $d_{r}$
  supplied by the plate cooling model of Equation~\eqref{eq:reference}, and
  the landward reach, which is retained for display and for locating the axis
  but enters no inversion. The landward and seaward window limits $x_{l}$ and
  $x_{s}$ of Equation~\eqref{eq:misfit} are indicated on the same axis, since
  they are the quantities the sweep of Section~\ref{sec:experiments} varies.
  Source: authors.}
  \label{fig:profilegeometry}
\end{figure}

Profiles are screened before use, under criteria applied identically at all
twenty margins. A profile is rejected if it crosses a seamount or fracture zone,
identified as a residual departure from a fifty-kilometre running median
exceeding five hundred metres over more than ten consecutive samples; if its
seaward reach intersects land or an adjacent trench; if fewer than half of its
nodes are flagged as measured in the type identifier grid \cite{GEBCO2023}; or
if no outer-rise crest can be located, taken as the shallowest sample between
twenty and two hundred kilometres seaward of the axis. The first criterion
matters most, because a seamount superimposes a short-wavelength load on the
bending signal and biases any fit that includes it; the third guards against
inverting a profile that is largely a gravity-derived prediction, not a
sounding. The criteria are applied without reference to the margin, and at every margin
of the inventory they remove a minority of profiles. They would remove all of
them at Cascadia, for the reason set out in Section~\ref{sec:studyarea}: a
trench buried under its own fan fill presents neither a depth minimum to serve
as the axis nor an outer rise to bound the fit. That is why the margin is not
in the inventory, and the outcome is reported rather than corrected, since
relaxing the criteria to admit it would mean fitting a bending profile to a
surface whose relief is depositional.

Two corrections convert the screened depth section into a deflection. The
sediment fill is removed isostatically, so that the surface being fitted is the
flexed igneous basement rather than the top of an accumulation whose thickness
varies along the profile \cite{Sykes1996}:
\begin{linenomath*}
\begin{equation}
  d_{c}(x) = d_{o}(x) + h_{s}(x)\,\frac{\rho_{m}-\rho_{s}}{\rho_{m}-\rho_{w}}
  \label{eq:sediment}
\end{equation}
\end{linenomath*}
where $d_{c}$ is the sediment-unloaded basement depth (m), $d_{o}$ the observed
depth (m), $h_{s}$ the total sediment thickness (m) from GlobSed
\cite{Straume2019}, and $\rho_{m}=3300$, $\rho_{s}=2300$ and
$\rho_{w}=1030$~kg~m$^{-3}$ the densities of mantle, sediment and sea water.

The second correction supplies the undeflected level from which bending is
measured. Oceanic lithosphere subsides as it cools, so the reference depth is a
function of plate age, not a constant, and it is taken from the GDH1
plate cooling model \cite{SteinStein1992}, the successor to the classical
half-space and plate fits of \citeA{ParsonsSclater1977}:
\begin{linenomath*}
\begin{equation}
  d_{r}(t) =
  \begin{cases}
    2600 + 365\,t^{1/2}, & t < 20\;\mathrm{Ma}\\[2pt]
    5651 - 2473\,\exp(-0.0278\,t), & t \geq 20\;\mathrm{Ma}
  \end{cases}
  \label{eq:reference}
\end{equation}
\end{linenomath*}
where $d_{r}$ is the reference depth (m) and $t$ the crustal age (Ma) taken from
the EarthByte grid \cite{Seton2020} as the segment mean. Age varies by less than
the grid resolution across a single profile at every margin treated here, so one
reference depth is applied per segment. The sampling is a single pass of
\texttt{grdtrack} over the grid of \citeA{Seton2020}, with the longitudes brought
onto the convention of that grid in the same command
(Listing~\ref{lst:agesample}); no age is interpolated across a node the model
leaves undefined, and a segment falling on such a node is carried forward with a
missing age rather than a filled one.

\begin{lstlisting}[caption={Sampling the EarthByte seafloor-age grid at the
mid-point of each segment. \texttt{-{}-GMT\_HISTORY=false} is not optional: a
\texttt{gmt.history} left by an earlier mapping command silently reapplies its
region and \texttt{grdtrack} then discards every point outside it, exiting
successfully with no output.},label={lst:agesample},language=bash]
# lon lat of each segment midpoint, wrapped from 0/360 onto -180/180
awk '{lon=$1; if (lon>180) lon-=360; print lon, $2}' segment_midpoints.txt |
    gmt grdtrack -Gage.2020.1.GTS2012.2m.nc -fg --GMT_HISTORY=false \
    > segment_ages.txt
# how many segments fall on an undefined node of the age model
awk '$3=="NaN" || $3==""' segment_ages.txt | wc -l
\end{lstlisting}

The deflection that both inversions fit is then the departure of the corrected
basement from that reference,
\begin{linenomath*}
\begin{equation}
  w(x) = d_{c}(x) - d_{r}
  \label{eq:deflection}
\end{equation}
\end{linenomath*}
with $w$ positive downward and $x$ the distance from the trench axis, positive
seaward, so that $w$ is large and positive at the axis and slightly negative
across the outer rise. Listing~\ref{lst:profiles}, in Appendix~\ref{app:code}, implements
Equations~\eqref{eq:sediment} to \eqref{eq:deflection} together with the
sampling and screening described above; the corrected profiles it returns are
the sole input to both inversions, which guarantees that the two branches see
identical data.

\subsection{Methodological Framework}\label{sec:framework}

Both inversions rest on the same physical model, so that any difference between
their results is attributable to the fitting procedure and not to the mechanics
assumed. The incoming plate is treated as a thin elastic plate of uniform
flexural rigidity, broken at the trench axis, resting on a fluid substratum and
loaded at its broken end by the plate boundary
\cite{Kirby2022Flexure,Caldwell1976}. Its deflection satisfies
\begin{linenomath*}
\begin{equation}
  D\,\frac{d^{4}w}{dx^{4}} + \Delta\rho\,g\,w = q(x)
  \label{eq:flexure}
\end{equation}
\end{linenomath*}
where $w$ is the deflection (m) of Equation~\eqref{eq:deflection}, $x$ the
distance from the axis (m), $D$ the flexural rigidity (N~m),
$\Delta\rho = \rho_{m}-\rho_{w} = 2270$~kg~m$^{-3}$ the density contrast across
the deflected surface, $g = 9.81$~m~s$^{-2}$ the gravitational acceleration, and
$q(x)$ the distributed vertical load (Pa). Seaward of the axis no surface load
acts on the sediment-unloaded basement, so $q(x)=0$ there and the bending is
driven entirely through the bending moment and shear force transmitted at the
broken end. In-plane forces are neglected, which keeps the recovered quantity a
single coefficient of a single ordinary differential equation; the consequences
are assessed in Section~\ref{sec:disc:limits}.

The rigidity carries the mechanical information sought, through the effective
elastic thickness \cite{Kirby2022Isostasy}:
\begin{linenomath*}
\begin{equation}
  D = \frac{E\,T_{e}^{3}}{12\left(1-\nu^{2}\right)}
  \label{eq:rigidity}
\end{equation}
\end{linenomath*}
where $E = 70$~GPa is Young's modulus, $\nu = 0.25$ Poisson's ratio and $T_{e}$
the effective elastic thickness (m). The cubic dependence is the reason
window-induced scatter matters: a ten per cent error in rigidity is a three per
cent error in thickness, but the converse amplification applies to any bias in
the fitted wavelength.

That wavelength is set by the flexural parameter,
\begin{linenomath*}
\begin{equation}
  \alpha = \left(\frac{4D}{\Delta\rho\,g}\right)^{1/4}
  \label{eq:alpha}
\end{equation}
\end{linenomath*}
in which $\alpha$ has units of length (m) and measures the distance over which
the plate distributes an applied load. It is the single quantity that fixes the
shape of the deflection: the position of the outer-rise crest and the decay
length of the swell are both proportional to $\alpha$, while the load fixes only
the amplitude. Recovering $T_{e}$ from a bending profile therefore amounts to
measuring a length scale.

Three further quantities follow from the same deflection and are recovered
alongside it:
\begin{linenomath*}
\begin{equation}
  M(x) = -D\,\frac{d^{2}w}{dx^{2}}, \qquad
  V(x) = -D\,\frac{d^{3}w}{dx^{3}}, \qquad
  \kappa(x) = \frac{d^{2}w}{dx^{2}}
  \label{eq:derived}
\end{equation}
\end{linenomath*}
where $M$ is the bending moment per unit length along strike (N), $V$ the shear
force per unit length (N~m$^{-1}$) and $\kappa$ the curvature (m$^{-1}$). All
three are derivatives of one function, so any method that returns a
differentiable deflection returns them without further fitting; this is the
property the physics-informed formulation exploits and the windowed baseline
does not.

\subsection{Baseline: Windowed Least-Squares Inversion}\label{sec:baseline}

The baseline reproduces established practice rather than a weakened version of
it, because the window experiment is intended to measure the procedure in
general use. For a broken plate loaded at its end by a vertical
shear force alone, Equation~\eqref{eq:flexure} with $q=0$ admits the closed-form
solution first applied to trenches by \citeA{Caldwell1976} and used, with minor
variants, in most subsequent profile studies
\cite{BodineWatts1979,Tandon2000,ContrerasReyesOsses2010}:
\begin{linenomath*}
\begin{equation}
  w(x) = \sqrt{2}\,w_{b}\,
  \exp\!\left[\frac{3\pi}{4}\left(1-\frac{x}{x_{b}}\right)\right]
  \cos\!\left(\frac{3\pi}{4}\,\frac{x}{x_{b}}\right)
  \label{eq:caldwell}
\end{equation}
\end{linenomath*}
where $w_{b}$ is the amplitude of the outer rise above the reference depth (m)
and $x_{b}$ the distance from the axis to the outer-rise crest (m), both marked
in Figure~\ref{fig:profilegeometry}. The two parameters separate cleanly:
$x_{b}$ carries the length scale, through $\alpha = 4x_{b}/3\pi$, and $w_{b}$
carries the magnitude of the end load. Inverting Equations~\eqref{eq:rigidity}
and \eqref{eq:alpha} gives the thickness directly,
\begin{linenomath*}
\begin{equation}
  T_{e} = \left[\frac{3\left(1-\nu^{2}\right)\Delta\rho\,g\,\alpha^{4}}{E}
  \right]^{1/3}
  \label{eq:te}
\end{equation}
\end{linenomath*}
so that a fitted crest distance maps to an elastic thickness without any
intermediate assumption. Equation~\eqref{eq:te} is used identically in both
branches of the study, which is what makes their outputs comparable.

The parameters are obtained by minimising the squared misfit between the
observed and predicted deflection over an interval of the profile:
\begin{linenomath*}
\begin{equation}
  \Phi\!\left(w_{b},x_{b}\right) =
  \frac{1}{N_{W}}\sum_{i\,:\,x_{l}\leq x_{i}\leq x_{s}}
  \left[w_{\mathrm{obs}}(x_{i}) - w(x_{i};w_{b},x_{b})\right]^{2}
  \label{eq:misfit}
\end{equation}
\end{linenomath*}
where $x_{l}$ and $x_{s}$ are the landward and seaward window limits (m),
$N_{W}$ the number of samples falling between them, and $w_{\mathrm{obs}}$ the
observed deflection. The two limits are the object of this study. They are not
estimated, not bounded by the data and not reported with the result: they are
constants supplied by the analyst before the minimisation begins, and every
value of $T_{e}$ in the profile-fitting literature is conditional on a
particular choice of them. Published choices differ in the seaward limit, in
whether the outer-rise crest is included or excluded, and in how far landward
the fit is allowed to approach the axis \cite{BryWhite2007,Manriquez2014}. The
minimisation itself is carried out by Levenberg--Marquardt least squares from a
grid of starting values, so that the recovered minimum is global within the
searched range and no part of the spread reported later is an artefact of
convergence.

\subsection{Proposed Method: Physics-Informed Neural Inversion}\label{sec:pinn}

The proposed formulation removes $x_{l}$ and $x_{s}$ from
Equation~\eqref{eq:misfit} rather than standardising them. In place of an
analytic solution fitted over a chosen interval, the deflection is represented
by a neural network whose derivatives are available in closed form through
automatic differentiation, and the governing equation is imposed as a residual
at points distributed across the whole profile
\cite{Raissi2019,Karniadakis2021}. The trial solution is
\begin{linenomath*}
\begin{equation}
  \hat{w}(x;\theta) = w_{0}\,
  \mathcal{N}\!\left(\frac{x}{L};\,\theta\right)
  \label{eq:trial}
\end{equation}
\end{linenomath*}
where $\mathcal{N}$ is a fully connected feed-forward network with hyperbolic
tangent activations, $\theta$ its weights and biases, $L$ the profile length (m)
and $w_{0}$ a fixed deflection scale (m) taken as the axial value of the
segment. Both scales are constants, not trainable quantities; they serve only to
present the network with inputs and targets of order unity, without which the
fourth derivative in the residual below is numerically ill-conditioned.

Because $\hat{w}$ is a composition of differentiable operations, its derivatives
with respect to $x$ are obtained by repeated application of the chain rule
through the computational graph, exactly and at machine precision, without
recourse to finite differences on a sampled profile. The physics residual is therefore
evaluable at any coordinate:
\begin{linenomath*}
\begin{equation}
  \mathcal{R}(x;\theta,T_{e}) = \frac{1}{\Delta\rho\,g\,w_{0}}
  \left[D(T_{e})\,\frac{d^{4}\hat{w}}{dx^{4}} + \Delta\rho\,g\,\hat{w}\right]
  \label{eq:residual}
\end{equation}
\end{linenomath*}
where $D(T_{e})$ is given by Equation~\eqref{eq:rigidity} and the leading factor
renders $\mathcal{R}$ dimensionless, so that residuals from segments of
differing amplitude are directly comparable and enter the loss on equal terms.
The elastic thickness appears in Equation~\eqref{eq:residual} as a trainable
coefficient of the same computational graph as $\theta$, and is recovered by the
same gradient descent; no separate inversion loop is required.

Training minimises a weighted sum of three terms,
\begin{linenomath*}
\begin{equation}
  \mathcal{L}(\theta,T_{e}) = \lambda_{d}\,\mathcal{L}_{\mathrm{data}}
  + \lambda_{p}\,\mathcal{L}_{\mathrm{phys}}
  + \lambda_{b}\,\mathcal{L}_{\mathrm{bc}}
  \label{eq:loss}
\end{equation}
\end{linenomath*}
whose components are
\begin{linenomath*}
\begin{equation}
  \mathcal{L}_{\mathrm{data}} = \frac{1}{N_{d}}\sum_{i=1}^{N_{d}}
  \left[\frac{\hat{w}(x_{i})-w_{\mathrm{obs}}(x_{i})}{w_{0}}\right]^{2},
  \qquad
  \mathcal{L}_{\mathrm{phys}} = \frac{1}{N_{c}}\sum_{j=1}^{N_{c}}
  \mathcal{R}(\xi_{j})^{2},
  \qquad
  \mathcal{L}_{\mathrm{bc}} =
  \left[\frac{\hat{w}(L)}{w_{0}}\right]^{2}
  + L^{2}\left[\frac{\hat{w}'(L)}{w_{0}}\right]^{2}
  \label{eq:lossterms}
\end{equation}
\end{linenomath*}
where $N_{d}$ is the number of sampled observations on the seaward reach,
$N_{c}$ the number of collocation points $\xi_{j}$ drawn from the same reach,
and $\lambda_{d}$, $\lambda_{p}$ and $\lambda_{b}$ dimensionless weights. The
boundary term enforces the only two conditions the problem supplies without an
analyst: that the deflection and its slope vanish in the far field. No condition
is imposed at the axis, since the bending moment and shear force there are
unknown; the axis enters solely as the origin of the coordinate, and it is
located by the data through the tracing of Section~\ref{sec:preprocessing}
rather than chosen. This is the precise sense in which the estimate is
window-free: the only interval boundary in the formulation is the seaward end of
the profile, which is fixed by the extraction geometry and identical at every
segment.

The three losses are minimised jointly across all segments of all twenty
margins, with a trunk network shared by every segment and one elastic thickness
per segment. Sharing is not a convenience. A single profile constrains a
fourth-order equation weakly, and an unshared network is free to reproduce the
observations with a deflection that satisfies the physics only loosely; a trunk
common to several hundred profiles must instead represent a family of shapes
that the governing equation admits, which acts as a regulariser that no
per-profile fit provides. Segment identity enters through a low-dimensional
learned embedding concatenated with the scaled coordinate, so that the shared
trunk can express margin-to-margin differences in shape while the recovered
thicknesses remain separate parameters and separate estimates.
Figure~\ref{fig:architecture} shows the forward and backward paths, and
Figure~\ref{fig:workflow} places the formulation within the whole procedure.
Listing~\ref{lst:loss}, in Appendix~\ref{app:code}, gives the implementation of
Equations~\eqref{eq:residual} to \eqref{eq:lossterms}.

\begin{figure}\centering
  \includegraphics[width=0.95\textwidth]{fig_architecture}
  \caption{Architecture of the physics-informed inversion. The scaled
  coordinate and the segment embedding enter a trunk network shared by every
  segment, whose output is the scaled deflection $\hat{w}$ of
  Equation~\eqref{eq:trial}. Four successive applications of automatic
  differentiation through the same graph supply the fourth derivative
  required by the residual of Equation~\eqref{eq:residual}, so that no finite
  difference is taken on a sampled profile. The elastic thickness enters that
  residual as a trainable coefficient of the same graph and is updated by the
  same gradient step as the network weights, which is why no separate
  inversion loop is required. Solid arrows, the forward pass; broken arrows,
  the backward pass. Source: authors.}
  \label{fig:architecture}
\end{figure}

\begin{figure}\centering
  \includegraphics[width=0.95\textwidth]{fig_workflow}
  \caption{Position of the physics-informed inversion within the whole
  procedure. The traced axes and the extracted profiles of
  Section~\ref{sec:preprocessing} are shared by both branches, so that the
  windowed baseline of Section~\ref{sec:baseline} and the window-free
  inversion of Section~\ref{sec:pinn} see identical corrected data and any
  difference between their results is attributable to the fitting procedure
  alone. The baseline branch carries the window limits as inputs and feeds
  the sweep of Section~\ref{sec:experiments}; the physics-informed branch
  carries none. Both converge on the comparison of
  Section~\ref{sec:res:validation}. Source: authors.}
  \label{fig:workflow}
\end{figure}

The training itself is driven by \texttt{train\_pinn.py}, whose invocation is
given in Listing~\ref{lst:run}. Four implementation choices in it are load
bearing and are recorded because they are not deducible from
Equations~\eqref{eq:trial} to \eqref{eq:lossterms} and because a run that
ignores them fails in ways that are not obvious from the loss curve. The whole
computation is carried out in double precision: the fourth derivative of
Equation~\eqref{eq:residual} spans some twenty-two orders of magnitude in SI
units over a profile five hundred and fifty kilometres long, and in single
precision the residual is dominated by rounding rather than by the physics, so
the loss falls while the recovered thickness drifts. The activation is the
hyperbolic tangent throughout, not for smoothness of the fitted deflection but
because the loss differentiates the network output four times, and a piecewise
linear activation has a second derivative that vanishes almost everywhere and
therefore yields an identically zero physics residual, which trains to a low
loss without ever having enforced the governing equation. The elastic thickness
is carried as $\log_{10}(T_{e}/\mathrm{km})$ rather than as $T_{e}$, so that it
cannot become negative under a gradient step and so that a step of fixed size is
a fixed relative change, the thickness being a scale rather than a location
parameter; it is bounded to $2 \leq T_{e} \leq 120$~km, the range of the
baseline of Section~\ref{sec:baseline} expressed as a thickness. The profile
length $L$ of Equation~\eqref{eq:trial} is one constant common to every segment,
fixed by the extraction geometry of Section~\ref{sec:preprocessing}, and the
deflection scale $w_{0}$ is per segment; both are constants of the graph and
neither is trained, and $w_{0}$ cancels identically in
Equation~\eqref{eq:residual}, so the choice affects conditioning only and not
the recovered thickness.

Four further choices concern the conditioning of the optimisation rather than
the formulation, and each was adopted after the alternative had been shown to
fail on the data of this study.

The scaled coordinate is not presented to the network directly but is lifted
into a set of random Fourier features, $\xi \mapsto [\sin(2\pi\mathbf{B}\xi),
\cos(2\pi\mathbf{B}\xi)]$, with the entries of $\mathbf{B}$ drawn once from a
zero-mean normal distribution of standard deviation $\sigma = 2$ and held fixed
thereafter. The motivation is the spectral bias of fully connected networks with
smooth activations, which fit the low-frequency content of a target long before
the high-frequency content \cite{Raissi2019,Karniadakis2021}. A bending profile
is not band-limited at low frequency: the trench wall and the outer-rise crest
together occupy the first fifth of a domain five hundred and fifty kilometres
long, while the remaining four fifths are nearly flat. The effect was quantified
by fitting the deflection of one Tonga segment with the data term alone, the
physics and boundary terms suppressed, for fifteen hundred iterations under
otherwise identical settings (Table~\ref{tab:encoding}). The encoding reduces
the residual by an order of magnitude, from forty to thirteen metres.
Table~\ref{tab:encoding} also records a negative result, since the
transformation concerned is ordinarily beneficial: recentring the coordinate to
$[-1,1]$, which improves conditioning for targets whose structure is distributed
across the domain, degrades the fit by a factor of twenty-five here, because it
places the sharply curved part of the profile at the edge of the input range
rather than at its origin. The bandwidth $\sigma$ is kept modest because the
residual of Equation~\eqref{eq:residual} differentiates the network output four
times, so that each retained frequency enters the physics term amplified by its
fourth power.

\begin{table}\centering
  \caption{Effect of the coordinate encoding on the achievable data misfit.
  One Tonga segment, the data term of Equation~\eqref{eq:lossterms} alone,
  fifteen hundred Adam iterations, with identical architecture, initialisation
  and learning rate in every row. The misfit is dimensionless, normalised by
  the deflection scale $w_{0}$; the residual is the same quantity expressed as
  a root-mean-square depth. Source: authors.}
  \label{tab:encoding}
  \begin{tabular}{lrr}
    \toprule
    Coordinate encoding & Data misfit & Residual (m) \\
    \midrule
    Plain coordinate, $x/L \in [0,1]$         & $7.4 \times 10^{-5}$ & 40 \\
    Recentred coordinate, $2x/L-1 \in [-1,1]$ & $1.9 \times 10^{-3}$ & 202 \\
    Random Fourier features, $\sigma = 2$     & $8.1 \times 10^{-6}$ & 13 \\
    \bottomrule
  \end{tabular}
\end{table}

The unit of estimation is the segment, and the deflection presented to the data
term is accordingly the pointwise median of the profiles of that segment on a
common one-kilometre grid, rather than the profiles themselves. Presenting them
individually imposes a lower bound on $\mathcal{L}_{\mathrm{data}}$ that no
network can pass, because a single curve $\hat{w}(x)$ cannot reproduce several
profiles that differ from one another; the bound is the mean square scatter of
the profiles about their own median, which for a Tonga segment is
$5.8 \times 10^{-4}$, of the same order as the misfit a converged fit should
attain. The optimiser would otherwise expend its capacity on along-strike
variability that the segment-wise formulation does not attempt to represent.
The median is preferred to the mean because a single profile crossing a
seamount or a fracture zone would displace the whole segment. The inter-profile
median absolute deviation is retained alongside the fitted residual, so that the
quality of a fit is judged against the scatter through which it passes rather
than against zero.

The second-stage refinement is evaluated on a fixed set of collocation points
and on every segment, whereas the first stage resamples both at each iteration.
The distinction is necessary rather than incidental. Stochastic resampling is
appropriate to a first-order method and acts as a regulariser, in that it
prevents the residual from being satisfied at a fixed and finite set of
coordinates; but a quasi-Newton method constructs its curvature estimate from
successive gradients of what it takes to be one fixed objective, and a line
search comparing function values across a changing objective terminates at its
first trial step. In an earlier configuration whose closure resampled, L-BFGS
executed fourteen of the two hundred iterations permitted before exiting. The
weights $\lambda_{d}$, $\lambda_{p}$ and $\lambda_{b}$ are frozen at the values
reached at the stage transition for the same reason.

Finally, the elastic thickness is assigned a learning rate distinct from, and
smaller than, that of the network parameters. The thickness is a single scale
parameter shared by every collocation point of a segment, whereas each network
weight contributes to the output through many paths, and the gradient
magnitudes of the two differ accordingly. Under a common rate the recovered
thickness of a Tonga segment executed a transient excursion to approximately
fifty-one kilometres before descending to a final value near twenty-two. The
excursion is not a blemish on Figure~\ref{fig:convergence} alone: at those
values the trial solution lies far from any deflection the plate equation admits
for the observed profile, the physics residual dominates the composite, and the
network is driven to compensate, so that iterations are expended undoing a
displacement that a smaller rate on the thickness prevents.

The dependency set is deliberately narrow. Training requires PyTorch, NumPy and
pandas and nothing else -- no plotting stack, no geospatial library -- so it is
installed in an environment of its own rather than beside the figure scripts.
The separation is a reproducibility measure and not a convenience: resolving
PyTorch into an environment that already carries the plotting and mapping stack
pulls those packages backwards, which changes the library versions recorded in
every figure of this paper. Training and rendering therefore run in different
environments, communicating only through the two comma-separated files the run
writes.

\begin{lstlisting}[caption={Environment and invocation for the joint training of
Section~\ref{sec:pinn}. The run writes a per-iteration history, the recovered
thickness of every segment, and a record of the seed, hyperparameters and
library versions sufficient to repeat it.},label={lst:run},language=bash]
# a dedicated environment: torch, numpy and pandas only, so that resolving
# PyTorch cannot downgrade the plotting stack used for the figures
conda create -n pinn python=3.11 pytorch numpy pandas -c pytorch
conda activate pinn

cd .../scripts/essential_data_scripts

# one margin first, to measure the iteration rate before committing
python3 train_pinn.py --profiles profiles_csv --outdir runs/test \
    --trenches tonga --adam 2000 --lbfgs 200

# the run reported here: the full inventory from one seed
python3 train_pinn.py --profiles profiles_csv --outdir runs/r02 \
    --adam 20000 --seed 0 --collocation 128

# outputs, in runs/r02:
#   history.csv    iteration, loss_total, loss_data, loss_phys, loss_bc, te_km
#                  components recorded UNWEIGHTED; loss_total is their plain
#                  sum, not the weighted objective the optimiser descends
#   segments.csv   segment, trench, te_km and the per-segment losses
#   lambda_log.csv the adaptive weights and gradient norms at each rebalancing
#   checkpoint.pt  trunk weights, segment embeddings and log10 T_e per segment
#   run.json       seed, hyperparameters, stopping state, library versions
\end{lstlisting}

\subsection{Experimental Design}\label{sec:experiments}

The window sweep is the primary experiment. For every retained segment the
baseline inversion of Equation~\eqref{eq:misfit} is repeated over a regular grid
of window limits: the landward limit $x_{l}$ takes eleven values from zero to
fifty kilometres seaward of the axis in five-kilometre steps, and the seaward
limit $x_{s}$ takes forty-eight values from eighty to five hundred and fifty
kilometres in ten-kilometre steps. The upper limit is the seaward reach of the
extracted profiles, so the grid spans every window the data admit. Of the five
hundred and twenty-eight combinations, four hundred and ninety-eight are
retained; the remainder are rejected by a single admissibility rule, that a
window must span at least eighty kilometres,
\begin{linenomath*}
\begin{equation}
  x_{s} - x_{l} \geq L_{\min}, \qquad L_{\min} = 80~\mathrm{km}
  \label{eq:minspan}
\end{equation}
\end{linenomath*}
because a shorter interval carries too few independent samples of the bending
wavelength to determine both $w_{b}$ and $x_{b}$ of Equation~\eqref{eq:misfit};
attempting it returns a formally converged fit whose flexural parameter is set
by the starting value rather than by the data. The rule is stated here because
it is the only constraint imposed on the sweep, and because it accounts for the
one systematic gap in the sensitivity surfaces described below. The grid is
chosen to contain the window conventions found in the published profile studies
rather than to bracket them symmetrically, so that the resulting spread is the
spread a reader of that literature would encounter. Every inversion in the sweep
uses the same profile, the same corrections and the same minimiser; the window
limits are the only quantity that varies. Algorithm~\ref{alg:sweep} states the
procedure.

Two further conditions are applied to the output of each inversion rather than
to its input, and both discard rather than repair. A fit whose flexural
parameter terminates on either bound of its admissible interval, $15 \leq \alpha
\leq 200$~km, has been stopped by the optimiser and not by the data, and is
recorded as missing; retaining such fits populates the surfaces with the two
thicknesses that correspond to the bounds themselves and inflates
$\Delta_{w}$ severalfold without any of it deriving from the bathymetry.
A fit whose root-mean-square residual exceeds $1.2$~km is likewise discarded,
the profile in that window being no longer describable by
Equation~\eqref{eq:deflection} to within the amplitude of the feature being
fitted. Both rejections are reported per segment with the sweep.

The behaviour of the baseline over the window plane is best seen before it is
reduced to a single number, and Figure~\ref{fig:windowmethod} shows it for four
margins. Each panel is the surface $T_{e}(x_{l}, x_{s})$ obtained by evaluating
the baseline at every admissible point of the grid on one profile, so the panel
is a direct portrait of the estimator: its horizontal axis is the seaward limit,
its vertical axis the landward limit, and its colour the thickness returned.
Read along the bottom edge, at a landward limit of zero, the panel reproduces the
one-dimensional sweep of Figure~\ref{fig:windowproblem}; the second axis extends
that sweep to the choice most studies leave unstated, namely how much of the
inner trench wall is admitted. The colour levels are spaced geometrically because the distribution of
recovered thickness is strongly right-skewed, and a linear spacing compresses
the plateau that occupies most of each panel into a single tint. The red
isoline drawn at the window-free estimate traces exactly the set of window
choices under which the classical method reproduces the window-free answer;
where it is absent, no admissible window reproduces that value. The blank wedge
at the left of every panel is not missing data but the region excluded by
Equation~\eqref{eq:minspan}, $x_{s} < x_{l} + 80$~km, whose boundary is the
diagonal running from $(80, 0)$ to $(130, 50)$, and the seaward axis begins at
$L_{\min}$ for the same reason. Three features recur on every margin and are
described here because they, and not any individual value, are what the
experiment is designed to expose. A narrow ridge of large recovered thickness
occupies the smallest seaward limits, where the window closes inside the
outer-rise crest and the fit compensates by lengthening the flexural wavelength.
A gradient in the opposite sense runs up the vertical axis, the recovered
thickness falling as more of the inner wall is excluded and the fitted section
is displaced seaward of the deflection maximum. Between them lies a broad
plateau, reached beyond roughly three hundred kilometres, on which the estimate
is stable to within a few per cent. That the plateau exists is not a defence of
the windowed method: the plateau of one margin does not coincide with the
plateau of another, none of the window conventions in the literature falls
inside all of them, and nothing internal to a single windowed inversion
identifies which part of its own surface is the plateau.

The four margins of Figure~\ref{fig:windowmethod} are not a random sample and
are not presented as one. They were selected purposively, before any surface was
computed, to span the three properties that control the shape of a bending
profile, so that a feature common to all four cannot be attributed to any one of
them. Plate age at the axis ranges from $40$ to $155$~Ma across the set --
Peru--Chile $40$~Ma, Aleutian $52$~Ma, Tonga $91$~Ma, Mariana $155$~Ma -- which
covers the thermal range of the whole inventory. Axial sediment thickness ranges
from $38$ to $597$~m, from the starved Tonga axis to the fed Aleutian, so the
correction of Section~\ref{sec:preprocessing} enters at very different
magnitudes. Axial depth ranges from $6.3$ to $10.3$~km, and the tectonic setting
spans intra-oceanic and continental margins and both orthogonal and strongly
oblique convergence. A single profile is used in each panel, not a segment
average, because the object being characterised is the estimator applied to one
section of bathymetry; averaging is introduced only in
Section~\ref{sec:res:window}, where the question changes from how the estimator
behaves to how much the margins differ.

The physics-informed neural network is trained once on all segments jointly.
Collocation points are drawn uniformly on the seaward reach and resampled at
every iteration, so that no fixed set of interior points can be fitted at the
expense of the rest of the profile. Optimisation proceeds in two stages: Adam
\cite{KingmaBa2015} with an initial learning rate of $3\times10^{-3}$ on the
network parameters and $10^{-3}$ on the elastic thickness, halved on decay,
followed by L-BFGS \cite{LiuNocedal1989} to convergence, a
sequence that is standard for this class of problem because first-order descent
locates the basin efficiently and the quasi-Newton stage resolves the residual
to a level finite-difference schemes cannot reach. The loss weights are set by
the gradient-norm balancing scheme of \citeA{Wang2021grad}, applied every one
hundred iterations, which prevents the data term from dominating the physics
term early in training. Training stops when the composite loss has not improved
by more than one part in $10^{4}$ over two thousand iterations.
Algorithm~\ref{alg:joint} gives the training loop, and
Table~\ref{tab:hyperparams} lists the searched ranges together with the retained
configuration.

All computation is in Python \texttt{3.12.13}. The networks and the automatic
differentiation are implemented in PyTorch \texttt{2.13.0}, run under Python
\texttt{3.11.15} in the separate training environment described below; array
handling uses NumPy \texttt{2.4.6} \cite{Harris2020}, tabulation uses pandas
\texttt{3.0.5},
and the baseline minimisation uses the Levenberg--Marquardt implementation of
SciPy \texttt{1.18.0} \cite{Virtanen2020}. Axis tracing, profile extraction and
all cartography use GMT \texttt{6.5.0} \cite{Wessel2019}, driven from shell
scripts rather than from a Python interface so that each map is a single
auditable command sequence, and diagnostic plots and schematics use Matplotlib
\texttt{3.11.1} \cite{Hunter2007}. Every run is seeded from a fixed value recorded with the code,
and the full environment specification, including operating system and pinned
library versions, is deposited alongside it so that the reported numbers can be
regenerated exactly.

\begin{figure}\centering
  \includegraphics[width=0.95\textwidth]{fig_ws_viridis}
  \caption{Behaviour of the windowed baseline over the plane of both window
  limits, for four margins spanning the range of the controlling variables.
  Each panel is the elastic thickness returned by the baseline inversion of
  Equation~\eqref{eq:misfit} applied to a single profile, evaluated at every
  admissible point of the sweep grid: the seaward limit $x_{s}$ horizontally and
  the landward limit $x_{l}$ vertically, both measured from the trench axis.
  \textbf{(a)} Mariana, plate age $155$~Ma, axial sediment $127$~m;
  \textbf{(b)} Peru--Chile, $40$~Ma, $161$~m; \textbf{(c)} Tonga, $91$~Ma,
  $38$~m; \textbf{(d)} Aleutian, $52$~Ma, $597$~m. Colour is the recovered
  thickness on a common scale across all four panels, with levels spaced
  geometrically so that equal colour steps correspond to equal ratios. White
  curves are isolines of recovered thickness, labelled in kilometres, and white
  circles the window choices adopted by previous profile studies. The red curve
  is the isoline at the window-free estimate returned by the physics-informed
  inversion for the same profile. The blank wedge at the left of every panel is
  the region excluded by Equation~\eqref{eq:minspan}. Panel annotations give
  $\Delta_{w}$ of Equation~\eqref{eq:windowmetric} for that profile. Software
  used for plotting figure: Python \texttt{3.12.13}, Matplotlib
  \texttt{3.11.1}, SciPy \texttt{1.18.0}, NumPy \texttt{2.4.6}. Data: GEBCO,
  GlobSed. Source: authors.}
  \label{fig:windowmethod}
\end{figure}

\begin{algorithm}
\caption{Window-sweep sensitivity experiment. The baseline inversion is repeated
over a regular grid of window limits and the recovered elastic thickness is
recorded for every combination.}\label{alg:sweep}
\begin{algorithmic}[1]
  \Require corrected profiles $\{x_i, w_i\}$ for every segment $s$
  \Require window grids $\mathcal{X}_l$, $\mathcal{X}_s$
  \Ensure sensitivity surface $T_e(s, x_l, x_s)$ and metric $\Delta_w(s)$
  \For{each segment $s$}
    \For{each $x_l \in \mathcal{X}_l$ and each $x_s \in \mathcal{X}_s$}
      \State select the samples with $x_l \leq x_i \leq x_s$
      \State minimise $\Phi(w_b, x_b)$ of Eq.~\eqref{eq:misfit} from a grid of
             starting values
      \State $\alpha \gets 4x_b/3\pi$; obtain $T_e$ from Eq.~\eqref{eq:te}
      \State store $T_e(s, x_l, x_s)$ and the misfit at the minimum
    \EndFor
    \State $\Delta_w(s) \gets$ Eq.~\eqref{eq:windowmetric} over the stored set
  \EndFor
\end{algorithmic}
\end{algorithm}

\begin{algorithm}
\caption{Joint physics-informed training across all segments, with a shared
trunk network and one elastic thickness per segment.}\label{alg:joint}
\begin{algorithmic}[1]
  \Require corrected profiles for all segments; scales $w_0(s)$, $L$
  \Ensure trained trunk $\theta$ and per-segment $\{T_e(s)\}$
  \State initialise $\theta$, the segment embeddings, and $T_e(s)$ at a common
         starting value
  \While{the stopping criterion is not met}
    \State draw a batch of segments and, for each, resample $N_c$ collocation
           points on the seaward reach
    \State evaluate $\hat{w}$ and, by automatic differentiation, its fourth
           derivative
    \State form $\mathcal{L}_{\mathrm{data}}$, $\mathcal{L}_{\mathrm{phys}}$ and
           $\mathcal{L}_{\mathrm{bc}}$ of Eq.~\eqref{eq:lossterms}
    \If{the iteration is a rebalancing step}
      \State update $\lambda_d$, $\lambda_p$, $\lambda_b$ from the gradient norms
    \EndIf
    \State update $\theta$ and $\{T_e(s)\}$ on $\nabla\mathcal{L}$
  \EndWhile
  \State refine with L-BFGS until convergence
  \State evaluate $M$, $V$ and $\kappa$ from Eq.~\eqref{eq:derived} on the
         trained solution
\end{algorithmic}
\end{algorithm}

\begin{table}\centering
  \caption{Hyper-parameters of the physics-informed inversion, with the range
  considered and the configuration retained. The ranges record the values
  examined while the method was developed; they were not swept systematically
  under the final configuration, and the consequence is entered in
  Table~\ref{tab:limitations}. The loss weights are set adaptively
  during training and their entry gives the initialisation. Where an iteration
  budget was not exhausted, the permitted and the executed count are both
  given. Source: authors.}\label{tab:hyperparams}
  \begin{tabular}{lll}
    \toprule
    Hyper-parameter & Search range & Retained \\
    \midrule
    Hidden layers & 3 to 8 & 5 \\
    Units per layer & 32 to 256 & 64 \\
    Activation & tanh, sine, softplus & tanh \\
    Fourier features, bandwidth $\sigma$ & 8 to 64, 1 to 8 & 16, 2.0 \\
    Embedding dimension & 2 to 16 & 4 \\
    Collocation points per segment & 128 to 2048 & 128 \\
    Segments per batch & 8 to all & 16 \\
    Learning rate, network & $10^{-4}$ to $10^{-2}$ & $3\times10^{-3}$ \\
    Learning rate, thickness & $10^{-4}$ to $10^{-2}$ & $10^{-3}$ \\
    Learning-rate decay factor & 0.1 to 0.9 & 0.5 \\
    Loss weights $\lambda_d:\lambda_p:\lambda_b$ & adaptive & $1:1:1$ initial \\
    Rebalancing interval (iterations) & 50 to 500 & 100 \\
    Adam iterations & $5\times10^{3}$ to $5\times10^{4}$ & $2\times10^{4}$ permitted, 4186 used \\
    L-BFGS iterations & 100 to 500 & 300 permitted, 316 recorded \\
    Patience, tolerance & --- & 2000 iterations, $10^{-4}$ \\
    Ensemble members & 5 to 50 & 1 (see Table~\ref{tab:limitations}) \\
    \bottomrule
  \end{tabular}
\end{table}

\subsection{Cartographic Mapping}\label{sec:mapping}

Every map in this study is produced by script from the open-access grids, with
no interactive editing, so that the cartography is reproducible on the same
terms as the inversion \cite{Wessel2019,Piskarev2019Lomonosov,Lemenkova202213}. The circum-Pacific
panels use an equidistant cylindrical projection centred on the date line, which
keeps the twenty margins on one sheet and preserves along-strike distance, the
quantity the segment bands encode. The base is shaded relief computed from the
bathymetric grid with illumination from the northwest at a fixed elevation, and
the shading is applied as an intensity layer beneath the data colours so that
relief and value remain separable to the eye.

Recovered parameters are rendered as coloured bands drawn along the traced axis,
one band per segment, at a width chosen so that adjacent segments remain
distinguishable at the reproduction size \cite{Lemenkova202017}. Discrete rather than continuous colour
steps are used for the elastic-thickness maps, because a reader comparing
margins needs to place a segment in a class, not to read an exact value from a
ramp. All colour scales are perceptually uniform and readable under the
common forms of colour-vision deficiency, and the diverging scales are anchored
on zero \cite{Crameri2020}; the sequential scales survive greyscale
reproduction, which is checked by converting each figure to luminance before it
is accepted.

The bands are produced by a join between two files that are stored separately,
because the geometry and the recovered values are computed at different stages
of the pipeline and share only the segment identity. The traced axes are held
as a GMT multisegment file with one record per trench, each headed by the
trench name and resampled to a common along-strike step, so that a single file
carries the geometry of all twenty margins (Listing~\ref{lst:axesfile}):
\begin{lstlisting}[caption={Structure of the traced-axis file: one
multisegment record per trench, headed by the trench name.},label={lst:axesfile},language=bash]
head -20 trench_axes_full.gmt   # > -L"Aleutian" then longitude latitude, 5 km
grep -c '^>' trench_axes_full.gmt   # 20: one axis record per trench
\end{lstlisting}
The recovered parameters are held per segment, several segments to a trench,
keyed by an identifier of the form \texttt{<trench>\_<nn>} that is numbered in
along-strike order. The number of segments varies by margin with the length of
the arc, and each segment carries an along-strike bin index in the \texttt{zone}
column, which confirms that the segments partition the axis into contiguous
equal-length bins rather than overlapping or leaving gaps
(Listing~\ref{lst:zones}):
\begin{lstlisting}[caption={Recovering the segment-to-trench mapping and the
along-strike bin index from the per-segment tables.},label={lst:zones},language=bash]
# segments per trench, in along-strike order (aleutian_01, _02, ...)
cut -d, -f1 segment_flexure.csv | sed 's/_[0-9]*$//' | sort | uniq -c
# the zone column gives each segment its along-axis bin
head -3 segment_zones.csv   # segment,trench,zone,w_b,Te,mdot,mmax,...
\end{lstlisting}
The merge therefore reads each trench axis, divides its node list into as many
contiguous equal runs as the trench has segments, and rewrites the header of
run $k$ to carry the recovered value of segment $k$ as a GMT $-Z$ attribute,
which \texttt{psxy} then maps to colour through the discrete palette. The output
is one multisegment file per mapped quantity, in which the geometry and the
value can no longer fall out of step because they are written together; the same
routine emits the elastic-thickness file and, by naming a different value
column, the bending-moment, shear, curvature and uncertainty files of
Figure~\ref{fig:parametermaps}. Longitudes are shifted to the nought-to-three-
hundred-and-sixty-degree convention in the same pass so that the bands cross the
date line without a seam. Naming a different value column is the whole of the
difference between the thickness map and its residual companion
(Listing~\ref{lst:resmap}), which is what allows the two sheets to be compared
band for band. Listing~\ref{lst:map} gives the rendering routine that
consumes the joined file.

\begin{lstlisting}[caption={Emitting a second mapped quantity from the same
join. Because both files inherit one partition of one axis file, band $k$ of the
residual map covers exactly the stretch of margin that band $k$ of the thickness
map covers.},label={lst:resmap},language=bash]
python3 build_te_segments.py --axes trench_axes_full.gmt \
    --values runs/r02/segments.csv --value te_km   --out te_segments.gmt
python3 build_te_segments.py --axes trench_axes_full.gmt \
    --values runs/r02/segments.csv --value loss_phys --out residual_segments.gmt
# the two files must agree band for band
grep -c '^>' te_segments.gmt residual_segments.gmt
\end{lstlisting}

\begin{lstlisting}[caption={The core of the routine rendering segment-wise
flexural parameters onto the circum-Pacific relief base map. The bands are
coloured by the per-segment value carried in each multisegment header, so the
geometry and the value cannot be separated.},label={lst:map},language=bash]
# shaded relief base, hinged palette, land masked
gmt grdcut "${RELIEF}" -R"${REG}" -Grelief.nc
gmt grdgradient relief.nc -A315 -Ne0.5 -Gshade.nc
gmt makecpt -Coleron -T-8000/6000 > relief.cpt

gmt begin fig_parametermaps pdf,png
  gmt subplot begin 2x2 -Fs8.1c/5.68c -M0.45c/1.35c -A"a)"+jTL \
                        -R"${REG}" -JQ8.1c
    gmt subplot set 0
      gmt grdimage relief.nc -Crelief.cpt -Ishade.nc -t30
      gmt coast -Ggray75 -Wthinnest,gray45 -Dl -A5000
      gmt basemap -Bxg60 -Byg30 -B+n          # graticule under the bands
      gmt plot m.gmt -Cm.cpt -W2.6p+cl        # +cl: colour by header -Z value
      gmt basemap -Bxa60f30 -Bya30f15 -BWeSN+t"Bending moment at the axis"
      gmt colorbar -Cm.cpt -DJBC+w6.6c/0.20c+h+o0/0.55c+e \
                   -Bx+l"M (10@+16@+ N)"
    # panels 1 to 3 repeat with v.gmt, k.gmt and d.gmt
  gmt subplot end
gmt end
\end{lstlisting}

\subsection{Validation and Performance Assessment}\label{sec:validation}

Validation proceeds on three tiers, each answering a different objection. The
first is recovery of a known answer: synthetic profiles are generated from
Equation~\eqref{eq:caldwell} for prescribed values of $T_{e}$ spanning the range
expected at the twenty margins, degraded with Gaussian noise at amplitudes
matched to the residual scatter of the observed profiles, and inverted by the
physics-informed formulation without disclosing the generating parameters.
Success is measured by the relative error in recovered thickness and by the root
mean square error between recovered and prescribed deflection. Because the
synthetic profiles are noise-free in their physics, this tier isolates the
performance of the estimator from any inadequacy of the elastic model.

The second tier is comparison with the windowed baseline on the observed
profiles. Both branches are applied to identical corrected profiles, and
agreement is assessed against the whole swept window population rather than
against a single privileged window, since the purpose is to establish where the
window-free estimate falls within the range the conventional method can produce.
The metrics are the root mean square error and mean absolute error of the fitted
deflection, the coefficient of determination, and the dimensionless physics
residual of Equation~\eqref{eq:residual}, which the baseline does not minimise
and which therefore measures how far each fitted profile departs from the
governing equation away from the fitted interval.

The third tier is comparison with independent published determinations at the
same margins, obtained by other data and other methods, including spectral
admittance and three-dimensional treatments. Agreement here cannot be exact,
because the published values are themselves conditional on window choices and on
differing physical assumptions, and a perfect match would be evidence of
circularity rather than of accuracy. The comparison is therefore read as a
consistency check on the order of magnitude and on the ranking of margins, and
it is reported with the span of the windowed baseline shown alongside, so that
a discrepancy can be judged against what the procedure alone can generate.
These metrics are reported in Section~\ref{sec:res:validation}.

\subsection{Statistical Analysis and Uncertainty Assessment}\label{sec:uncertainty}

Uncertainty in the recovered parameters is estimated by ensemble. The joint
training is repeated with independently initialised weights, independently drawn
collocation sequences and independently shuffled segment batches, and the
dispersion of the recovered thickness across ensemble members is taken as its
uncertainty. This captures the components that dominate for this class of
estimator, namely sensitivity to initialisation and to the stochastic sampling
of the physics constraint, without the cost of full posterior inference
\cite{Psaros2023}. The derived quantities of Equation~\eqref{eq:derived} are
evaluated separately within each member and summarised by the same statistics,
so that their uncertainties inherit the correlation structure of the ensemble
instead of being propagated through an assumed linearisation. Uncertainty
arising from the input grids is treated separately, by repeating the inversion
with the sediment and reference-depth corrections perturbed within the
tolerances quoted by their providers.

The window dependence of the baseline is quantified by a single metric, defined
here because it is a contribution of this study rather than a reported result.
For each segment $s$, over the set $\mathcal{W}$ of swept windows,
\begin{linenomath*}
\begin{equation}
  \Delta_{w}(s) = \frac{\max_{W\in\mathcal{W}} T_{e}(s,W)
  - \min_{W\in\mathcal{W}} T_{e}(s,W)}
  {\operatorname{median}_{W\in\mathcal{W}} T_{e}(s,W)}
  \label{eq:windowmetric}
\end{equation}
\end{linenomath*}
where $T_{e}(s,W)$ is the thickness recovered from segment $s$ under window $W$.
The metric is dimensionless, so segments of differing strength are directly
comparable, and it is normalised by the median, not the mean, so that a
few windows that place the fit across the outer-rise crest cannot inflate the
denominator. Its interpretation is exact: $\Delta_{w}$ is the fraction of the
typical recovered thickness that is attributable to nothing but the analyst's
choice of interval, evaluated on a fixed profile with fixed physics. Geological
variation cannot enter it, because every term in Equation~\eqref{eq:windowmetric}
is computed from the same segment. The interquartile range of the same
population is reported alongside it as a robust counterpart, insensitive to the
extremes of the swept grid.

% =========================================================================
\section{Results}\label{sec:results}

\subsection{Profile Inventory and Data Quality}\label{sec:res:quality}
The extraction of Section~\ref{sec:preprocessing} yields one thousand eight
hundred and sixty-seven profiles that pass every screening criterion, cast at
five-kilometre spacing along the traced axes and each sampled at one-kilometre
intervals from one hundred and fifty kilometres landward of the axis to five
hundred and fifty kilometres seaward. The largest samples come from the longest
arcs -- one hundred and eighty-seven profiles at the Aleutian, one hundred and
seventy-nine at the Mariana, one hundred and sixty-three at Middle America and
one hundred and fifty at the Philippine -- and the smallest from the shortest,
with fourteen at Palau and six at Hikurangi; the per-margin counts are given in
Table~\ref{tab:trenches}. Grouped into the along-strike segments of the same
table, the sample supports one hundred and thirty-seven segments that carry a
flexural estimate, fewer than the partition of the axis lengths admits, because
a segment left empty by screening is dropped rather than filled. Of the retained
profiles, one thousand seven hundred and four exhibit a resolvable outer-rise
crest and one hundred and sixty-three do not, the latter concentrated on the
sediment-fed margins where the rise is onlapped; both classes are retained,
since the inversion of Section~\ref{sec:pinn} requires no feature to be picked.
Plate age is unavailable at one hundred and seventy-one profiles, all of them
within the polygons the age model leaves undefined near continental margins and
young spreading centres, and those profiles are excluded from the age
relation of Section~\ref{sec:res:parameters} but from nothing else.
(Table~\ref{tab:trenches}). The retained sample covers every margin of the
inventory. Cascadia is not among them: its trench is buried beneath its fan
fill and yields no traceable axis, and it is therefore excluded from the
inventory itself rather than carried through it as an empty row, as set out in
Sections~\ref{sec:studyarea} and \ref{sec:preprocessing}.

The inversion reported below is trained on a subset of that inventory, and the
reduction from one to the other is read directly off the two tables rather than
asserted (Listing~\ref{lst:inventory}). Of the
segments carrying a flexural estimate, seventy-one enter the run, drawn from
twenty margins, because the per-segment profile cap of
Section~\ref{sec:preprocessing} admits at most eight profiles to the stack and
segments whose screened profiles fall below the minimum are not stacked at all.
Five hundred and fifty-three profiles are stacked in total, sixty-five of the
seventy-one segments receiving the full complement of eight and the remainder
between three and seven (Listing~\ref{lst:stacks}). Stacking yields 38{,}627 median-profile samples on the
common one-kilometre grid, a median of 544 nodes per segment over the
550~km seaward reach. The deflection scale $w_{0}$ of the stacked profiles has
a median of 2712~m across the segments and spans 694 to 6440~m; the shortfall
below the cap falls on the margins whose traced axis is shortest.

\begin{lstlisting}[caption={Reading the segment inventory off the two
per-segment tables. The first count is the full set carrying a flexural
estimate; the second is the subset that enters the neural
inversion.},label={lst:inventory},language=bash]
tail -n +2 segment_flexure.csv | wc -l       # 137 segments with an estimate
tail -n +2 runs/r02/segments.csv | wc -l     # 71 segments in the run
grep -c '^>' te_segments.gmt                 # 71 mapped bands, one per segment
cut -d, -f2 runs/r02/segments.csv | tail -n +2 | sort -u | wc -l   # 20 margins
# segments per margin, most-populated first
cut -d, -f2 runs/r02/segments.csv | tail -n +2 | sort | uniq -c | sort -rn
\end{lstlisting}

\begin{lstlisting}[caption={Profiles stacked per segment, and the node count of
the stacked deflections. The cap of eight profiles per segment is visible as
the mode of the first distribution.},label={lst:stacks},language=bash]
# distribution of profiles per stacked segment
cut -d, -f3 runs/r02/segments.csv | tail -n +2 | sort -n | uniq -c
# total stacked profiles, and median-profile samples on the 1 km grid
awk -F, 'NR>1{s+=$3} END{print s, "profiles stacked"}' runs/r02/segments.csv
tail -n +2 runs/r02/deflection.csv | wc -l   # 38627 samples
\end{lstlisting}

\subsection{Magnitude of the Window Dependence}\label{sec:res:window}
The baseline of Section~\ref{sec:baseline}, implemented in
Listing~\ref{lst:baseline}, was first evaluated at a single
window of the kind an analyst supplies before the minimisation begins, the
landward limit at the axis and the seaward limit at 250~km, and the resulting
fits are shown for one section of each margin in
Figure~\ref{fig:baselinefits}. The fitted curves follow the
observed sections closely: the residual remains within about one kilometre over
the fitted interval on the majority of the twenty panels, and departs from that
band only where a feature the elastic model does not represent crosses the
section, at the seamount on the Japan section near 100~km, over the seaward half
of the Palau and Yap sections, and on the San Cristobal section landward of the
outer-rise crest. The thicknesses recovered from these fits nevertheless span
4~km at Hikurangi and Puysegur to 113~km at Manila, a factor of twenty-eight
about a median of 14~km, with fifteen of the twenty margins below 20~km and the
remaining five --- Ryukyu, Mariana, Yap, Vanuatu and Manila --- between 39 and
113~km. Compared against the window-free column of
Table~\ref{tab:windowspread}, and allowing that the panels carry one section
rather than the margin as a whole, fourteen of the twenty agree to within a
factor of two, those same five exceed the window-free estimate by factors of
three to seven, and Hikurangi falls an order of magnitude below it. Goodness of
fit therefore does not separate the two groups, and nothing internal to a single
windowed inversion distinguishes a thickness of 13~km from one of 113~km.

\begin{figure}\centering
  	\includegraphics[width=0.98\textwidth]{fig_baselinefits}
  \caption{Classical windowed least-squares fits to one bathymetric section from
  each of the twenty circum-Pacific margins, all at the same window: the
  landward limit at the trench axis and the seaward limit at 250~km. One panel
  per margin, labelled with the margin name and with the effective elastic
  thickness $T_{e}$ recovered from that section. Within a panel the grey line is
  the observed section, plotted as depth below sea level and positive downward;
  the broken vertical line marks the trench axis at zero distance; the open
  circle marks the outer-rise crest; and the red line is the analytical
  deflection of Equation~\eqref{eq:caldwell} fitted by minimising
  Equation~\eqref{eq:misfit}, drawn over the fitted window alone and returned to
  the depth axis by restoring the reference depth of
  Equation~\eqref{eq:reference}. The narrow strip beneath each panel gives the
  residual between observation and fit in blue, on a vertical scale common to
  all twenty panels so that misfit can be compared across margins. Distance is
  measured from the trench axis, positive seaward. Software used for plotting
  figure: Python \texttt{3.12.13}, Matplotlib \texttt{3.11.1}, SciPy
  \texttt{1.18.0}, NumPy \texttt{2.4.6}. Data: GEBCO bathymetry. Source:
  authors.}
  \label{fig:baselinefits}
\end{figure}

The sweep of Section~\ref{sec:experiments} was run on every retained segment of
all twenty margins, and Figure~\ref{fig:windowsensitivity} presents the
resulting sensitivity surfaces, one panel per trench. Each panel aggregates the
individual profile surfaces of that margin by taking, at every point of the
window grid, the median of the thicknesses recovered from the sampled profiles;
the panel is therefore a property of the margin rather than of any profile
chosen to represent it, and the arbitrariness of that choice, which
Figure~\ref{fig:windowmethod} is subject to by construction, is removed. Colour
is the recovered thickness normalised by the median of the same panel, so that
margins whose absolute thickness differs by a factor of five remain comparable
as sensitivity surfaces and the colour axis carries the quantity that
Equation~\eqref{eq:windowmetric} measures. The scale is geometric rather than
linear, which places a doubling and a halving of the thickness equally far from
the neutral colour. The three features identified in
Section~\ref{sec:experiments} are present on every margin, which establishes
that the window dependence is a property of the estimator and not of any
particular section of seafloor.

Three conventions govern the panels. The blank wedge at the left of each is the
region excluded by Equation~\eqref{eq:minspan}, within which no window spans
the eighty kilometres the fit requires; blank regions elsewhere, at Palau,
Manila, Middle America and Hikurangi, are windows at which no sampled profile
of that margin returned an admissible fit, and are therefore an absence of data
rather than of colour. The surfaces are resampled by bicubic spline for
legibility, and every number reported below is computed on the unsmoothed grid.
The annotated $\Delta_{w}$ is evaluated on the median surface and is for that
reason systematically smaller than the per-profile values of
Table~\ref{tab:windowspread}, since taking the median across profiles before
taking the range suppresses the variation the metric is intended to capture;
the two are not interchangeable, and the tabulated per-profile population
is the one that supports the aggregate reported below.

Swept over the full grid of window limits, the recovered thickness at a single
margin varies by a factor that is large everywhere and extreme on the least
constrained arcs. Measured per profile, the window-induced spread
$\Delta_{w}$ of Equation~\eqref{eq:windowmetric} ranges from $1.34$ at
Kuril--Kamchatka to $7.26$ at Vityaz, with a median of $3.02$ across the twenty
margins; that is, the median margin returns a largest windowed thickness some
three times its smallest, and at the extreme the ratio exceeds seven. Seventeen
of the twenty margins exceed a factor of two and ten exceed a factor of three,
so a thickness quoted without its window is uncertain by more than the
between-margin differences the parameter is used to interpret. The values are
tabulated by margin in Table~\ref{tab:windowspread} and shown as sensitivity
surfaces in Figure~\ref{fig:windowsensitivity}; the surfaces carry the
per-window-median $\Delta_{w}$, which for the reason given above is
systematically smaller than the per-profile figure, spanning $0.42$ at Mariana
to $6.39$ at Vityaz with a median of $1.16$
(Figure~\ref{fig:windowsensitivity}, Table~\ref{tab:windowspread}).

The spread is greatest where the sample is smallest, so it must be read against
the length of the arc before it is taken as a property of the margin alone. The
five largest values fall on the shortest and most thinly sampled trenches ---
Vityaz ($\Delta_{w}=7.26$, three profiles available), Palau ($5.75$, fourteen),
Manila ($4.39$, fifty-six), Puysegur ($4.38$, twenty-two) and Philippine
($3.88$) --- and across the twenty margins $\Delta_{w}$ correlates negatively
with the number of profiles available ($r=-0.55$). Part of the spread at the
short arcs is therefore sampling noise rather than genuine window sensitivity,
and the aggregate figure above is correspondingly weighted towards the
well-sampled majority, at which $\Delta_{w}$ nonetheless remains above two.

The dependence is also visible directly, without aggregating into a surface, by
sweeping one limit at a time and watching the population of recovered
thicknesses move with it (Figure~\ref{fig:windowspread}). Advancing the seaward
limit from 80 to 540~km lifts the median recovered thickness across all 428
profiles from 10.6 to 15.3~km, a drift of 45~per cent produced by nothing but
the choice of where to stop fitting, and the interquartile spread among
profiles widens over the same interval from 76 to 155~per cent of the median.
The landward limit acts over a far shorter lever with a comparable effect, the
median rising from 13.6 to 17.1~km as the limit retreats 50~km from the axis.
Margins differ in how strongly they respond: at Vityaz the spread stays between
25 and 31~per cent of the median at every setting, whereas at Puysegur the
median itself swings between 5.9 and 13.3~km and the spread reaches 209~per
cent. In every panel the interval-free estimate of Section~\ref{sec:pinn} lies
below the windowed population at all settings, a systematic offset returned to
in Section~\ref{sec:disc:limits}.

\begin{figure}\centering
  \includegraphics[width=0.98\textwidth]{fig_windowspread}
  \caption{The window dependence swept one limit at a time, from 72{,}452
  windowed inversions of 428 profiles at twenty margins. Panels (a) and (b)
  pool every margin and sweep the seaward and the landward limit respectively;
  panels (c) and (d) sweep the seaward limit at the least and the most
  window-sensitive margin, selected by the smallest and largest interquartile
  spread over windows whose seaward limit is at least 200~km rather than by
  inspection. In every panel the purple curve is the median recovered
  thickness over the profiles contributing at that window setting, on the left
  axis, and the shaded band its interquartile range; the orange broken curve is
  that interquartile range expressed as a percentage of the median at the same
  setting, on the right axis, and is the quantity
  Equation~\eqref{eq:windowmetric} formalises. The broken vertical line marks a
  conventional seaward limit of 250~km, or the axis for the landward sweep. The
  dotted horizontal line is the interval-free estimate of
  Section~\ref{sec:pinn}, which has no window and therefore no position on the
  abscissa. Software used for plotting figure: Python \texttt{3.12.13},
  Matplotlib \texttt{3.11.1}, NumPy. Source: authors.}
  \label{fig:windowspread}
\end{figure}

The same sweep resolved by margin rather than by window limit shows how
unevenly the sensitivity is distributed across the belt
(Figure~\ref{fig:trenchbox}). Each box is the interquartile range of the
windowed thickness at one margin, pooled over every admissible window with a
seaward limit of at least 200~km, so the height of the box is the window
sensitivity of that margin and the line within it the value an analyst would
most often report. The margins are ordered by median, and the spread of box
heights is large: the interquartile range is a median of 88~per cent of the
margin's own thickness, but it falls to 29~per cent at Vityaz, where the
windowed estimate is nearly stable, and rises to 187~per cent at Puysegur,
where it is barely defined. The interval-free estimate lies at or below the
lower quartile of almost every box, the same offset seen in the pooled sweep.

The sensitivity is not organised geographically. Grouping the margins into
northern, north-western, western, south-western and eastern sectors of the
basin leaves the median sensitivity between 72 and 105~per cent in every
sector, with the tall and the short boxes intermixed within each: the
south-western group alone contains both the least sensitive margin, Vityaz, and
some of the most sensitive, Puysegur and New Britain. Nor does the sensitivity
track the recovered thickness or the length of the arc; across the twenty
margins the interquartile spread is uncorrelated with the median thickness
($r = 0.00$) and only weakly related to the number of windows available
($r = 0.14$). What the box height does follow is the morphology of the
individual margin, the width and clarity of its outer rise and the presence or
absence of features that a fixed window includes or excludes as its limits
move, which is a property of each trench's bathymetry rather than of its
position in the basin or its tectonic setting. The window sensitivity is
therefore an artefact of the estimator meeting a particular seafloor shape, not
a signal with a regional structure, which is the reason it cannot be reasoned
away by restricting the study to one sector or one class of margin.

\begin{figure}\centering
  \includegraphics[width=0.98\textwidth]{fig_trenchbox}
  \caption{Window sensitivity of the classical estimator, one box per margin,
  from the same sweep as Figure~\ref{fig:windowspread}. For each margin the box
  spans the interquartile range of the effective elastic thickness recovered by
  windowed least squares over every admissible window with a seaward limit of
  at least 200~km, the line within the box is the median, and the whiskers span
  the 5th to 95th percentiles; outliers beyond the whiskers are omitted. The
  height of a box is thus the window sensitivity of that margin: a tall box is a
  margin whose recovered thickness depends strongly on the fitting interval, a
  short box one that is robust to it. Margins are ordered by median thickness
  and coloured across a spectral scale for identification. The dashed line is
  the interval-free estimate of Section~\ref{sec:pinn}, which has no window and
  lies at or below the lower quartile of almost every margin. These are windowed
  \emph{baseline} thicknesses, characterising the classical estimator, not the
  recovered thicknesses of the neural inversion, which are tabulated per margin
  in Table~\ref{tab:tebytrench}. Software used for plotting figure: Python
  \texttt{3.12.13}, Matplotlib \texttt{3.11.1}, NumPy. Source: authors.}
  \label{fig:trenchbox}
\end{figure}

\begin{figure}\centering
  \includegraphics[width=0.98\textwidth]{fig_windowsensitivity}
  \caption{Dependence of the recovered effective elastic thickness on both
  window limits, for each of the twenty circum-Pacific margins, panels (a) to
  (t). Axes are as in Figure~\ref{fig:windowmethod}: the seaward window limit
  $x_{s}$ horizontally and the landward limit $x_{l}$ vertically, both in
  kilometres from the trench axis. Each panel is the per-window median over the
  profiles sampled at even spacing along that margin, the sample size being
  given with each panel as $n$ of the profiles available. Colour is the
  recovered thickness divided by the median of the same panel, on a geometric
  scale spanning one octave either side of unity; windows falling outside that
  octave are shown at the ends of the scale. Isolines are drawn at fixed ratios
  and labelled. Blank areas carry no admissible fit. Panel annotations give
  $\Delta_{w}$ of Equation~\eqref{eq:windowmetric} evaluated on the median
  surface, together with the number of profiles entering that surface and the
  number available at the margin; every profile that returns an admissible fit
  is used, so the surfaces involve no subsampling. Software used for plotting figure: Python \texttt{3.12.13},
  Matplotlib \texttt{3.11.1}, SciPy \texttt{1.18.0}, NumPy \texttt{2.4.6}.
  Data: GEBCO, GlobSed. Source: authors.}
  \label{fig:windowsensitivity}
\end{figure}

\begin{table}\centering
  \caption{Window-induced spread in recovered effective elastic thickness, by
  trench. The minimum, maximum and interquartile range of $T_{e}$ (km) are
  taken over the pooled admissible windowed fits of the sampled profiles; the
  minimum and maximum approach the bounds of the admissible interval on most
  margins, so the interquartile range is the informative measure of scale. The
  metric $\Delta_{w}$ of Equation~\eqref{eq:windowmetric} is the per-profile
  median, dimensionless. The final column is the window-free physics-informed
  estimate (km) for comparison. Source: authors.}
  \label{tab:windowspread}
  \begin{tabular}{lrrrrr}
    \toprule
    Trench & $T_e$ min & $T_e$ max & $T_e$ IQR & $\Delta_{w}$ & Window-free \\
    \midrule
    Aleutian         &    4 &  111 &    8 &  1.88 & 9 \\
    Kuril--Kamchatka &    4 &  111 &    6 &  1.34 & 8 \\
    Japan            &    4 &  109 &   10 &  2.16 & 10 \\
    Izu--Bonin       &    4 &  110 &    8 &  2.13 & 11 \\
    Mariana          &    4 &  111 &   20 &  2.87 & 10 \\
    Yap              &    4 &  106 &   11 &  2.36 & 9 \\
    Palau            &    4 &  111 &    9 &  5.75 & 9 \\
    Ryukyu           &    4 &  111 &   33 &  2.58 & 6 \\
    Manila           &    4 &  105 &    6 &  4.39 & 27 \\
    Philippine       &    4 &  111 &   15 &  3.88 & 9 \\
    New Britain      &    4 &  111 &   45 &  1.87 & 23 \\
    San Cristobal    &    4 &  111 &   29 &  3.17 & 10 \\
    Vityaz           &    4 &  105 &    4 &  7.26 & 9 \\
    Vanuatu          &    4 &  111 &   43 &  2.11 & 21 \\
    Tonga            &    4 &  110 &   13 &  3.36 & 10 \\
    Kermadec         &    4 &  107 &   17 &  2.56 & 10 \\
    Hikurangi        &    4 &   58 &    9 &  3.73 & 67 \\
    Puysegur         &    4 &  110 &   13 &  4.38 & 5 \\
    Middle America   &    4 &  111 &    8 &  3.21 & 8 \\
    Peru--Chile      &    4 &  108 &   12 &  3.67 & 9 \\
    \bottomrule
  \end{tabular}
\end{table}

\subsection{Convergence and Fitted Bending Profiles}\label{sec:res:fits}
The configuration of Table~\ref{tab:hyperparams} carries 19{,}428 trainable
parameters: 19{,}073 in the trunk, 284 in the segment embeddings and one
thickness for each of the seventy-one segments, the sixteen Fourier frequencies
being drawn once and held fixed. The joint training was run once over all
seventy-one segments from the fixed seed on a single processor, and completed
4502 iterations in 13{,}809~s, three hours and fifty minutes, at an average of
3.07~s per iteration. The first stage
terminated at iteration 4186 of the twenty thousand permitted, on the criterion
of Section~\ref{sec:experiments}: the composite loss had not improved by more
than one part in $10^{4}$ over the preceding two thousand iterations. Its
running median over two hundred iterations reached $3.3\times10^{-2}$ by
iteration 2000 and stood at $3.8\times10^{-2}$ at the transition, so the budget
consumed after the plateau is the patience window itself. The second stage
recorded a further 316 iterations before exiting on its own convergence test.
The training loop itself is given in Listing~\ref{lst:train} and the ensemble
procedure in Listing~\ref{lst:ensemble}, both in Appendix~\ref{app:code}.
Both stage boundaries and the final loss split are read from the iteration log
rather than from the console output of the run (Listing~\ref{lst:history}).

\begin{lstlisting}[caption={Reading the convergence record. The stage switch and
the stopping criterion are recorded in \texttt{run.json}; the per-iteration
budget comes from the elapsed time divided by the row count of the
history.},label={lst:history},language=bash]
head -1 runs/r02/history.csv    # iteration,loss_total,loss_data,loss_phys,loss_bc,te_km
tail -1 runs/r02/history.csv    # final iteration and its loss split
tail -n +2 runs/r02/history.csv | wc -l          # 4502 iterations
grep -E 'stage_switch|stop_reason|elapsed_s' runs/r02/run.json
\end{lstlisting}

The three loss terms fall by unequal amounts and at different stages
(Figure~\ref{fig:convergence}a). Over the whole run the data term falls from
$8.3\times10^{-2}$ to $2.0\times10^{-2}$, a factor of 4.2; the physics term from
$2.2\times10^{-2}$ to $1.5\times10^{-2}$; and the boundary term from
$7.7\times10^{-3}$ to $3.3\times10^{-5}$, a factor of 234. The quasi-Newton stage
accounts for a disproportionate share of the last two: across its 316
iterations the physics term falls by 19.7 per cent and the boundary term by
83.2 per cent, while the data term changes by 0.33 per cent. The refinement
therefore acts on satisfaction of the governing equation and of the seaward
conditions rather than on the fit to the bathymetry, which the first stage has
already resolved. Single-iteration values of the composite carry the sampling
noise of the first stage, whose batch is sixteen segments and whose collocation
points are redrawn at every step; over the final thousand first-stage
iterations the composite has a relative standard deviation of 32 per cent about
its own mean, and its smallest recorded value, $1.07\times10^{-2}$ at iteration
2185, is one such fluctuation and not a minimum the run departs from. The
composite recorded in the history is the unweighted sum of the three terms and
is not the objective descended, which carries the adaptive weights.

Those weights separate the terms by more than an order of magnitude
(Table~\ref{tab:lambdas}). Gradient-norm balancing raises $\lambda_{d}$ from
1.01 at the first rebalancing step to 3.29 at the last and $\lambda_{b}$ from
1.96 to 46.33, the latter reaching its imposed ceiling of 50 at five of the
forty-one steps, while $\lambda_{p}$ remains at unity to within
$8.8\times10^{-4}$ throughout. The values frozen at the stage transition are
$\lambda_{d} = 3.295$, $\lambda_{p} = 1.000$ and $\lambda_{b} = 46.334$.

The elastic thickness of the tracked segment descends monotonically from the
common initial value of 25.0~km, through 14.4~km at iteration 1000 and 8.2~km
at iteration 3000, to 7.06~km at the stage transition, and settles at 7.37~km
(Figure~\ref{fig:convergence}b). Over the final hundred iterations it varies
between 7.10 and 7.37~km, a range of 3.7 per cent of the converged value.

The recovered deflection reproduces the stacked profiles to a median
root-mean-square residual of 126~m, with an interquartile range of 84 to 228~m,
against a median deflection scale of 2712~m; the normalised data misfit of
Equation~\eqref{eq:lossterms} is the square of that ratio and has a median of
$3.0\times10^{-3}$. Forty-seven of the seventy-one segments fall below 150~m and
twenty-five below 100~m, the best fit being 40~m at an Aleutian segment and the
worst 749~m at New Britain. The residual is smaller than the inter-profile
median absolute deviation on thirty-seven of the seventy-one segments, and the
two are equal at the median, their ratio being 0.99, so that the fits reach the
scatter of the profiles about their own median rather than passing below it
(Figure~\ref{fig:profilefits}). The same twenty margins under the windowed
baseline at its 250~km window are shown in Figure~\ref{fig:baselinefits}.

\begin{table}\centering
  \caption{Adaptive loss weights of Equation~\eqref{eq:loss} at the first and
  last rebalancing step of the first stage, and the values frozen at the stage
  transition and held through the quasi-Newton refinement. Weights are
  dimensionless. The interval between steps is one hundred iterations, giving
  forty-one steps over the 4186 iterations of the first stage. Source:
  authors.}
  \label{tab:lambdas}
  \begin{tabular}{lrrr}
    \toprule
    Weight & Step 1 (iter. 100) & Step 41 (iter. 4100) & Frozen \\
    \midrule
    $\lambda_{d}$, data     &  1.01 &  3.29 &  3.295 \\
    $\lambda_{p}$, physics  &  1.00 &  1.00 &  1.000 \\
    $\lambda_{b}$, boundary &  1.96 & 46.33 & 46.334 \\
    \bottomrule
  \end{tabular}
\end{table}

\begin{figure}\centering
  \includegraphics[width=0.90\textwidth]{fig_convergence}
  \caption{Convergence of the joint physics-informed training over all
  seventy-one segments. Both panels share the iteration axis, and the broken
  vertical line marks the transition from Adam to L-BFGS at iteration 4186,
  labelled above the frame. \textbf{(a)} the composite loss $\mathcal{L}$ of
  Equation~\eqref{eq:loss}, solid black, and its three components as recorded,
  before weighting: the data term $\mathcal{L}_{\mathrm{data}}$ broken, the
  physics term $\mathcal{L}_{\mathrm{phys}}$ dash-dotted and the boundary term
  $\mathcal{L}_{\mathrm{bc}}$ dotted, on a logarithmic scale; all four are
  dimensionless. \textbf{(b)} the elastic thickness recovered for one Aleutian
  segment, initialised at 25~km in common with every segment, with the
  horizontal broken line at the converged value of 7.37~km. Line style
  distinguishes the components independently of colour, so both panels survive
  greyscale reproduction. Software used for plotting figure: Python
  \texttt{3.12.13}, Matplotlib \texttt{3.11.1}, NumPy \texttt{2.4.6}. Source:
  authors.}
  \label{fig:convergence}
\end{figure}

\begin{figure}\centering
  	\includegraphics[width=0.98\textwidth]{fig_profilefits}
  \caption{Observed and recovered bending profiles for one representative
  segment of each of the twenty trenches, one panel per trench and each
  labelled with the trench name. Within a panel, grey circles are the corrected
  deflection of Equation~\eqref{eq:deflection}, sampled at one kilometre; the
  solid line is the deflection recovered by the physics-informed inversion; the
  open circle marks the outer-rise crest. The narrow strip beneath each panel
  shows the residual between the two on a common vertical scale, so misfit can
  be compared across trenches. Depth is positive downward and distance is
  measured from the trench axis, positive seaward. Software used for plotting
  figure: Python \texttt{3.12.13}, Matplotlib \texttt{3.11.1}. Data: GEBCO
  bathymetry, GlobSed sediment thickness, EarthByte seafloor age. Source:
  authors.}
  \label{fig:profilefits}
\end{figure}

\subsection{Recovered Flexural Parameters}\label{sec:res:parameters}
Across the seventy-one segments the window-free effective elastic thickness has
a median of 9.4~km and an interquartile range of 8.1 to 10.7~km, within a full
range of 4.0 to 67.0~km (Figure~\ref{fig:temap},
Table~\ref{tab:tebytrench}). Half of the population therefore occupies a band
2.7~km wide, and the outer decades of the range are carried by few segments:
the fifth and ninety-fifth percentiles are 5.4 and 26.0~km. The smallest value
is 4.0~km at a Kuril-Kamchatka segment and the largest 67.0~km at the single
Hikurangi segment, with Puysegur, Yap and Middle America also returning values
below 5~km. The quartiles and the extremes quoted here are computed from the
recovered column of the run table, not from the rounded values of
Table~\ref{tab:tebytrench} (Listing~\ref{lst:tedist}).

\begin{lstlisting}[caption={The recovered thickness distribution, computed from
the run table. Sorting numerically before taking the quantiles matters: the
column is right-skewed and a lexical sort would place 9.4 after
67.0.},label={lst:tedist},language=bash]
# median, quartiles and extremes of te_km over the 71 segments
cut -d, -f6 runs/r02/segments.csv | tail -n +2 | sort -g |
    awk '{v[NR]=$1} END{printf "min %.1f  Q1 %.1f  med %.1f  Q3 %.1f  max %.1f\n",
          v[1], v[int(NR*0.25)], v[int(NR*0.5)], v[int(NR*0.75)], v[NR]}'
# margins whose median exceeds 20 km
awk -F, 'NR>1{s[$2]=s[$2]" "$6} END{for(t in s) print t, s[t]}' \
    runs/r02/segments.csv | sort
\end{lstlisting}

Between margins the medians span a factor of fourteen, from 4.8~km at Puysegur
to 67.0~km at Hikurangi. Four margins carry a single segment and therefore have
no along-strike range, their tabulated entries repeating. The margins represented by four or more segments
are confined to a much narrower interval, 7.5~km at Kuril-Kamchatka to 20.6~km
at Vanuatu, and eleven of the twenty have medians between 7.5 and 10.6~km.
Within a margin the along-strike spread is smallest where the sample is largest
and the arc is uniform: Japan varies by 0.4~km over three segments and
Peru-Chile by 1.3~km over five, whereas Mariana spans 6.6 to 19.6~km over six
and Kuril-Kamchatka 4.0 to 13.1~km over five.

The largest thicknesses coincide with the poorest fits. The six segments
exceeding 20~km carry a median root-mean-square residual of 472~m against 113~m
for the remaining sixty-five, and the rank correlation between recovered
thickness and residual over all segments is 0.31. Restricting the population to
the forty-seven segments whose residual falls below 150~m, which span thirteen
margins, leaves the median at 9.3~km and the interquartile range at 7.9 to
10.1~km, and reduces the full range to 4.0 to 18.0~km: the upper extreme of the
unrestricted distribution is contributed entirely by segments the elastic model
fits poorly.

Over the 47 segments retained by the residual screen the bending moment at the
axis has a median of $6.1\times10^{15}$~N (interquartile range
$2.5\times10^{15}$ to $1.3\times10^{16}$), the shear force a median of
$8.5\times10^{11}$~N\,m$^{-1}$ ($6.1\times10^{11}$ to $1.4\times10^{12}$), and
the maximum curvature along the profile a median of
$1.6\times10^{-6}$~m$^{-1}$ ($1.1\times10^{-6}$ to $2.2\times10^{-6}$). The
curvature maximum falls within 31~km of the axis on every segment, at the
trench wall where the plate is most sharply bent, so the extremum is a property
of the recovered bending and not of the seaward end of the domain. All three
follow from the trained solution by Equation~\eqref{eq:derived} in a single
evaluation pass, without refitting

\begin{figure}\centering
  	\includegraphics[width=0.98\textwidth]{fig_temap}
  \caption{Window-free effective elastic thickness of the subducting plate at
  the twenty circum-Pacific trenches. Each 148 kilometre segment is drawn as a
  coloured band along the traced axis, on shaded bathymetric relief, with plate
  boundaries in broken black. The colour scale is discrete, perceptually
  uniform and colour-blind safe, so that a segment is placed in a class rather
  than read from a continuous ramp, and it survives greyscale reproduction.
  Segments rejected by the profile screening carry no band. Software used for
  plotting figure: GMT \texttt{6.5.0}. Data: GEBCO
  bathymetry, PB2002 plate boundaries. Source: authors.}
  \label{fig:temap}
\end{figure}

\begin{table}\centering
  \caption{Window-free effective elastic thickness by trench, from the joint
  inversion of Section~\ref{sec:pinn}. Segments is the number entering the
  inversion; median, minimum and maximum are taken over those segments, and the
  residual is the median root-mean-square misfit to the stacked profile.
  Thicknesses in kilometres, residuals in metres. Source: authors.}
  \label{tab:tebytrench}
  \begin{tabular}{lrrrrr}
    \toprule
    Trench & Segments & Median & Min & Max & Residual \\
    \midrule
    \multicolumn{6}{l}{\textit{Northwestern and western Pacific}} \\
    Aleutian         &  7 &   8.8 &   7.4 &  18.0 &  79 \\
    Kuril-Kamchatka  &  5 &   7.5 &   4.0 &  13.1 & 104 \\
    Japan            &  3 &  10.1 &   9.8 &  10.2 &  98 \\
    Izu-Bonin        &  4 &  10.6 &   9.6 &  11.8 & 103 \\
    Mariana          &  6 &   9.8 &   6.6 &  19.6 & 156 \\
    Yap              &  3 &   8.6 &   4.4 &  11.9 & 234 \\
    Palau            &  1 &   8.6 &   8.6 &   8.6 & 144 \\
    Ryukyu           &  3 &   6.4 &   5.9 &   7.7 & 126 \\
    Manila           &  2 &  27.4 &   6.6 &  48.2 & 456 \\
    Philippine       &  5 &   8.7 &   7.8 &  11.2 & 109 \\
    \midrule
    \multicolumn{6}{l}{\textit{Southwestern Pacific}} \\
    New Britain      &  3 &  22.6 &  18.4 &  29.8 & 507 \\
    San Cristobal    &  2 &   9.9 &   8.7 &  11.1 & 230 \\
    Vanuatu          &  4 &  20.6 &  10.2 &  27.6 & 334 \\
    Vityaz           &  1 &   8.8 &   8.8 &   8.8 & 415 \\
    Tonga            &  4 &  10.3 &   9.1 &  10.7 & 106 \\
    Kermadec         &  5 &  10.1 &   7.4 &  12.7 & 102 \\
    Hikurangi        &  1 &  67.0 &  67.0 &  67.0 & 329 \\
    Puysegur         &  1 &   4.8 &   4.8 &   4.8 & 106 \\
    \midrule
    \multicolumn{6}{l}{\textit{Eastern Pacific}} \\
    Middle America   &  6 &   7.8 &   4.9 &  10.9 &  74 \\
    Peru-Chile       &  5 &   9.3 &   8.6 &  10.0 &  96 \\
    \midrule
    \textit{All segments} & 71 &   9.4 &   4.0 &  67.0 & 126 \\
    \bottomrule
  \end{tabular}
\end{table}

\begin{figure}\centering
  \includegraphics[width=0.98\textwidth]{fig_parametermaps}
  \caption{The remaining recovered flexural parameters, over the same
  circum-Pacific base as Figure~\ref{fig:temap} and drawn as coloured
  along-strike bands so that the four fields can be compared segment by
  segment. \textbf{(a)} bending moment at the axis, in units of $10^{16}$~N;
  \textbf{(b)} shear force at the axis, in units of $10^{12}$~N~m$^{-1}$;
  \textbf{(c)} maximum curvature along the profile, in units of
  $10^{-7}$~m$^{-1}$; \textbf{(d)} the deflection scale at the axis, $w_{b}$,
  in metres, which is the amplitude the other three carry and is shown so that
  their along-strike pattern can be read against it. Panels (a) to (c) follow
  Equation~\eqref{eq:derived} and are evaluated on the same differentiable
  solution as the thickness, over all seventy-one segments of the run. Each
  panel carries its own scale, the quantities having different dimensions and
  different signs: the shear is negative at every segment and the moment and
  curvature change sign along strike, so a band at the pale end of a scale is
  not a small value but a value near the limit stated on that bar. Bathymetry
  is shaded relief from the GEBCO\_2023 Grid. Software used for plotting
  figure: GMT \texttt{6.5.0}. Source: authors.}
  \label{fig:parametermaps}
\end{figure}

\subsection{Validation and the Windowed Baseline}\label{sec:res:validation}
On the observed profiles the inversion attains a median root-mean-square
deflection error of 126~m and a median normalised data misfit of
$3.0\times10^{-3}$, against a median physics residual of $5.8\times10^{-3}$ and a
median boundary residual of $9.0\times10^{-8}$ (Table~\ref{tab:metrics}). The
boundary term is thus satisfied to five orders of magnitude below the two terms
that compete, and the residual budget at convergence is divided between the fit
to the bathymetry and the governing equation in the ratio one to two. The
recovered solution accounts for a median 92 per cent of the variance of the
stacked deflection within a segment, with an interquartile range of 87 to 95
per cent; the weakest segment retains 48 per cent and no segment falls below
the variance of its own mean, so the fit is nowhere worse than uninformative.

The recovered thicknesses may also be read against the windowed baseline of
Section~\ref{sec:baseline} applied to the same profiles. Over the windows an
analyst would defend, with a seaward limit of at least 200~km, the baseline
returns a median thickness of 15~km with an interquartile range of 8 to 27~km
across the swept grid, against the interval-free median of 9.4~km recovered
here. The two distributions overlap but the interval-free estimate sits at the
lower edge of the windowed population, an offset examined in
Section~\ref{sec:disc:limits}; what the comparison establishes is not that one
value is correct but that a single windowed number, quoted without its window,
carries the 40~per cent spread of Figure~\ref{fig:windowspread}, whereas the
interval-free estimate carries none. A comparison against independent published
determinations is deferred: a like-for-like table would require each source to
be reduced to the same margins and the same definition of the fitted quantity,
which is beyond the scope of this methodological study.

\begin{table}\centering
  \caption{Fit metrics for the physics-informed inversion on the observed
  profiles, with the recovered and the windowed-baseline thickness for
  comparison. Deflection errors in metres, the coefficient of determination and
  the residuals dimensionless, thicknesses in kilometres. Each fit metric is
  the median over the 71 segments, the coefficient of determination computed per
  segment against that segment's own mean deflection before the median is taken.
  The baseline thickness is the median and interquartile range of the windowed
  least-squares estimate over the swept grid with a seaward limit of at least
  200~km. Source: authors.}
  \label{tab:metrics}
  \begin{tabular}{lr}
    \toprule
    Metric & Observed profiles \\
    \midrule
    Deflection error, median (m)      & 126 \\
    Deflection error, IQR (m)         & 84--228 \\
    Normalised data misfit, median    & $3.0\times10^{-3}$ \\
    Physics residual, median          & $5.8\times10^{-3}$ \\
    Boundary residual, median         & $9.0\times10^{-8}$ \\
    Coefficient of determination      & 0.92 \\
    \midrule
    Recovered thickness, median (km)  & 9.4 \\
    Baseline thickness, median (km)   & 15 \\
    Baseline thickness, IQR (km)      & 8--27 \\
    \bottomrule
  \end{tabular}
\end{table}

\subsection{Reduction of the Recovered Thickness Against Bending}\label{sec:res:yielding}
The recovered thickness is an effective quantity: bending at a trench drives the
outer fibres of the plate past the yield envelope, so the section that continues
to behave elastically is thinner than the mechanical plate
\cite{GoetzeEvans1979,WattsBurov2003}. That reduction is the quantity of
interest, and the joint formulation permits it to be tested against its supposed
cause without appeal to an external model. The curvature is recovered from the
same differentiable solution as the thickness (Equation~\eqref{eq:derived}), so
the degree of reduction and the bending that is supposed to produce it are
measured on one profile, at one place, by one inversion.

The reduction is expressed as the ratio of the recovered thickness to the depth
of the 600~\textdegree C isotherm at that segment's plate age, the isotherm
conventionally taken as the base of the mechanical lithosphere. A ratio near
unity is a plate bending elastically through its full mechanical thickness; a
ratio well below unity is a plate that has yielded through part of its section.
Figure~\ref{fig:yielding} plots that ratio against the maximum curvature
recovered along each segment. Both quantities come from the run of
Section~\ref{sec:pinn} and neither is imported: the isotherm depth follows from
the age of Section~\ref{sec:datamaterials}, and the curvature from
Equation~\eqref{eq:derived}. Where the elastic description fits the observed
bending poorly the ratio is not interpretable as a degree of yielding, and the
margins at which that happens are mapped in Figure~\ref{fig:residualmap}.

Across the 47 retained segments the recovered thickness is a median of 21~per
cent of the depth to the 600~\textdegree C isotherm, with an interquartile
range of 18 to 30~per cent. The reduction is therefore large and, within this
sample, nearly uniform: regressing the ratio on the logarithm of the curvature
gives a slope of $-0.23$ per decade, in the expected sense that more sharply
bent segments retain a smaller elastic fraction, but the relation accounts for
8~per cent of the scatter and a rank correlation of $-0.26$ does not reach
significance at this sample size ($p = 0.08$). The evidence is thus for a
pervasive reduction rather than one modulated segment by segment by the local
bending, which is consistent with the near-uniformity of the recovered
thickness itself (Section~\ref{sec:res:parameters}). Three segments return a ratio
above unity, which is not physically admissible for a purely thermal plate;
these are Hikurangi, Manila and San Cristobal, all among the poorly fitted
segments the residual screen removes, so the non-physical values arise where
the elastic description fails rather than where it holds.

\begin{figure}\centering
  	\includegraphics[width=0.92\textwidth]{fig_yielding}
  \caption{Reduction of the recovered elastic thickness against the bending
  that produces it. Vertical axis, the ratio of the window-free thickness to the
  depth of the 600~\textdegree C isotherm at that segment's plate age, so unity,
  marked by the dashed line, is a plate bending elastically through its full
  mechanical thickness. Horizontal axis, the maximum curvature along the
  segment on a logarithmic scale, recovered from the same solution by
  Equation~\eqref{eq:derived}. One marker per segment, coloured by trench; open
  markers are segments whose root-mean-square deflection error exceeds 150~m,
  drawn for completeness but excluded from every statistic, as in
  Figure~\ref{fig:teage}. The solid blue line is the median ratio over the
  retained segments and the blue band its interquartile range; these summarise
  the population and are the primary result, since a near-constant ratio is
  what the figure shows. The purple diamonds are the median ratio in four
  equal-count bins of curvature, a guide to any monotone drift rather than a
  tested relation. The red line is the least-squares fit on $\log_{10}\kappa$
  and the red band its 95~per cent confidence interval; the band straddles the
  horizontal median across the whole range, which is the visual statement that
  the slope is not distinguishable from zero. The axis is capped just above
  unity to fill the panel with the populated range; three segments lie above
  the cap and are marked by arrows at the top with their values and names,
  Hikurangi, Manila and San Cristobal, all poorly fitted, drawn hollow, and
  excluded from the fit and the median. A ratio above unity is not physically
  admissible for a thermal plate, so these arise where the elastic description
  fails rather than where it holds. Both quantities are recovered by the
  inversion of Section~\ref{sec:pinn}, so the reduction is regressed on the
  bending without recourse to an external thermal or rheological model.
  Software used for plotting figure: Python \texttt{3.12.13}, Matplotlib
  \texttt{3.11.1}, SciPy, NumPy. Source: authors.}
  \label{fig:yielding}
\end{figure}

\begin{figure}\centering
  	\includegraphics[width=0.95\textwidth]{fig_residualmap}
  \caption{Spatial distribution of the physics residual of
  Equation~\eqref{eq:residual} after training, drawn segment by segment over
  the same circum-Pacific base as Figure~\ref{fig:temap}. The residual is
  dimensionless, so segments of differing bending amplitude are directly
  comparable. High values identify margins at which a plate of uniform rigidity
  cannot reproduce the observed bending, and are therefore a diagnostic of
  model inadequacy and not of estimator failure; the windowed baseline returns
  no comparable quantity. Software used for plotting figure: GMT
  \texttt{6.5.0}. Source: authors.}
  \label{fig:residualmap}
\end{figure}

% =========================================================================
\section{Discussion}\label{sec:discussion}

\subsection{How Much of the Published Disagreement Is Methodological}\label{sec:disc:window}

The window dependence measured here is not scatter but drift, and its direction
follows from how least squares distributes weight. Every sample inside the
interval enters the misfit of Equation~\eqref{eq:misfit} with the same weight,
while the information about the flexural parameter is concentrated in one part
of the profile: the position of the outer-rise crest, which fixes $x_{b}$ and
therefore $\alpha$. A window that stops short of the crest removes that
constraint and leaves only the steep near-axis limb, where amplitude and
wavelength trade off against each other, so the minimisation compensates a
shortened wavelength with a larger end load and returns a stiffer plate. A
window extended far beyond the crest has the opposite defect: it fills the sum
of squares with samples whose deflection is close to zero and whose departures
from it are seafloor roughness, not bending, so the fit is pulled by
structure the elastic model was never meant to reproduce. Between those two
regimes the estimate passes through no stationary point, which is why sweeping
the limit produces a systematic trend and not a plateau with noise on it.

A second mechanism compounds the first. The governing equation constrains the
fourth derivative of the deflection, whereas the windowed misfit constrains only
the deflection itself. Two solutions that agree closely on $w$ within a chosen
interval can differ substantially in curvature, and it is curvature that carries
the rigidity. The classical procedure therefore has no mechanism for detecting
that a good-looking fit is mechanically wrong, and the analyst receives no
signal from the misfit that the window has misled it. This is the reason the
physics residual of Equation~\eqref{eq:residual} is reported for the baseline as
well as for the proposed method: it is the diagnostic the baseline lacks.

The consequence for the literature is a reallocation of variance rather than a
dismissal of it. If two-fifths of the interquartile spread among published determinations at a
single margin can be reproduced by moving the window on one profile, then the
physical explanations advanced for that spread have been asked to account for
variance they did not generate. Inelastic yielding, bend-faulting, hydration and
along-strike variability remain real and remain necessary
\cite{ContrerasReyes2007,Craig2014,KaoChen1996,MoscosoContrerasReyes2012}; what
changes is the size of the residual disagreement they are required to explain,
and therefore the confidence with which any one of them can be preferred over
another. The point applies with particular force to compilations, where values
obtained under different conventions are placed on a common axis and the
convention is not recorded \cite{BryWhite2007,Manriquez2014}.

The metric $\Delta_{w}$ is not, however, a property of the method alone. A
segment with a well-developed outer rise constrains $x_{b}$ from a wide range of
windows and is correspondingly insensitive to the choice; a segment whose rise
is subdued, buried beneath trench fill or overprinted by a subducting seamount
offers the crest to only a narrow band of windows and is correspondingly
fragile. Window dependence is therefore expected to be largest exactly where the
flexural signal is weakest, which is also where a physical interpretation of a
low recovered thickness is most tempting. That coincidence is the practical
danger the sweep exposes: the margins at which the procedure is least
trustworthy are not flagged as such by the misfit it returns.

\subsection{Physical Interpretation of the Window-Free Parameters}\label{sec:disc:physics}

Once the window is removed, the recovered thickness can be read against plate
age without the confounding of procedure. The expectation from a cooling plate
is that the effective elastic thickness tracks the depth to a controlling
isotherm and therefore grows approximately as the square root of age
\cite{SteinStein1992,ParsonsSclater1977,CalmantCazenave1987}, and the
circum-Pacific gravity inversion of \citeA{HunterWatts2016} places that isotherm
at $390 \pm 10$~\textdegree C. Figure~\ref{fig:teage} sets the window-free
estimates against that expectation, and they do not follow it: the recovered
thickness is near uniform with plate age rather than rising as its square root
(Section~\ref{sec:res:parameters}), so no single isotherm is resolved and none
is quoted here. Were an age dependence present, the implied temperature rather
than the thickness would be the quantity to compare across studies, because it
is insensitive to the absolute calibration of $T_{e}$ in a way the thickness is
not: a systematic bias in recovered thickness shifts the isotherm by a much
smaller amount through the
square-root relation.

\begin{figure}\centering
  \includegraphics[width=0.92\textwidth]{fig_teage}
  \caption{Window-free effective elastic thickness against the age of the
  subducting plate. One marker per segment, coloured by trench and keyed in the
  legend, with vertical bars giving the ensemble standard deviation. Broken
  lines, the depth to isotherms of a plate cooling model at the temperatures
  labelled on each curve, which a plate deforming elastically to a
  temperature-controlled depth would follow. Solid line and shaded band, the
  relation fitted to the segments here with its confidence interval. Thickness
  in kilometres and age in millions of years. Software used for plotting
  figure: Python \texttt{3.12.13}, Matplotlib \texttt{3.11.1}. Source:
  authors.}
  \label{fig:teage}
\end{figure}

The recovered thickness is an effective quantity and sits below the mechanical
thickness of the plate, because bending at a trench drives the outer fibres past
the yield envelope and the plate ceases to behave elastically through its full
depth \cite{GoetzeEvans1979,WattsBurov2003}. The size of that reduction is not a
nuisance but a measurement: it is the degree to which the plate has failed under
its own bending. What the joint formulation adds is the means to test the
reduction against its supposed cause. A single-parameter fit returns $T_{e}$ and
nothing else, so the inference that yielding explains a low value rests on
comparison with an external thermal model. Recovering the curvature
$\kappa$ from the same differentiable solution
(Equation~\eqref{eq:derived}) allows the reduction to be regressed directly on
the bending that is supposed to produce it, segment by segment, with both
quantities measured on the same profile. That regression is
Figure~\ref{fig:yielding} and Section~\ref{sec:res:yielding}.

The bending moment and shear force carry information of a different sort. They
are the loads transmitted at the broken end of the plate, and their along-strike
variation is a map of what the plate boundary applies to the incoming
lithosphere, not of how the lithosphere resists it. Segments where the
moment is large but the thickness is unremarkable indicate strong loading of a
normal plate; segments where the thickness falls while the moment does not
indicate weakening. Separating those two cases is not possible from $T_{e}$
alone, and it is the reason the parameter set is recovered jointly rather than
sequentially. Departures below the age trend are expected where the incoming
plate is pervasively bend-faulted and hydrated, since both processes reduce
strength without changing age \cite{ContrerasReyes2007,Craig2014}, and the
segments concerned should also carry elevated outer-rise normal faulting.

\subsection{Comparison With Classical and Machine-Learning Inversion}\label{sec:disc:comparison}

Against the windowed baseline the proposed formulation is expected to agree
where the classical procedure is best conditioned and to depart where it is
not: agreement at segments with a broad, well-expressed outer rise, and
divergence at segments where no window both reaches the crest and avoids the
far-field roughness. The baseline comparison of
Section~\ref{sec:res:validation} is therefore read as a test of conditioning
rather than of correctness. A window-free estimate that falls inside the swept
span is not thereby confirmed;
it is shown to be one of the values the classical method could have produced,
and the sweep shows how much choosing was required to produce it.

Differences from published circum-Pacific compilations have
identifiable sources. \citeA{HunterWatts2016} invert free-air gravity in place
of bathymetry, allow $T_{e}$ to vary across the outer rise, and
require profiles at least 750~km long; each of those three choices should raise
their recovered thickness relative to a uniform-plate bathymetric fit on a
shorter profile, and the first also removes the sediment correction as a source
of error. \citeA{BryWhite2007} reappraised the trench compilation and attributed
much of its scatter to modelling choices, a conclusion this study supports on
different evidence and extends by isolating one of those choices and measuring
it. The older profile determinations
\cite{Caldwell1976,BodineWatts1979,Tandon2000} used the same broken-plate
solution reproduced in the baseline here, so any difference from them is
attributable to data vintage and to the window, which is precisely what the
sweep quantifies.

The relation to data-driven inversion is one of kind, not of degree. Neural
inversions of seismic impedance and petrophysical properties learn a mapping
from observations to parameters out of a training set of examples
\cite{Das2019,ChenSchuster2020,Zhao2024}, and their central difficulty is
whether that mapping survives outside the distribution it was trained on. No
such training set exists for flexural parameters, because the true elastic
thickness has never been measured anywhere; the quantity is defined by the model
that recovers it. The formulation used here therefore never learns an inverse
operator. The network is a differentiable representation of one deflection
profile, the governing equation supplies the constraint that a training set
would otherwise have to supply, and generalisation in the usual sense does not
arise. The cost of that immunity is the absence of amortisation: every new
margin requires an optimisation, whereas a trained operator would require a
forward pass.

Within physics-informed inversion the flexural problem is unusually
well-conditioned, and much of the literature in this area addresses far harder
cases. Reviews of the field, and its applications
to wave-equation and rock-physics inversion, contend with many coupled unknown
fields and with residuals that constrain those fields only weakly
\cite{Schuster2024,Li2024Probabilistic,Li2024Bayesian,Zhang2025GINN,ChengAlkhalifah2023,Qu2022Separate}.
Flexure presents a single ordinary differential equation with one physically meaningful
coefficient per segment, so the recovered coefficient can be interpreted
directly instead of being treated as a by-product of fitting a field. That is
what makes it a suitable case for asking whether the physics constraint can
replace an analyst's interval, and the answer obtained here would not transfer
automatically to a problem with a less identifiable parameter.

\subsection{Transferability of the Formulation}\label{sec:disc:transfer}

The ingredient that makes the window disappear is not the neural network. It is
the pairing of a governing equation expressible as a differentiable residual
with a parameter that enters that equation as a coefficient. Given both, the
constraint can be imposed at arbitrary coordinates instead of at sampled data
points, and the estimate stops depending on which samples the analyst chose to
include. Any inverse problem posed on a profile and governed by an ordinary
differential equation of that form is a candidate for the same PINN
treatment: thermal subsidence and
heat-flow profiles across cooling lithosphere, flexure under seamount, fracture
zone and mountain loads \cite{Sandwell1982,StewartWatts1997,Bodine1981},
compaction profiles through sedimentary sections. Where no governing equation is
available, the residual can be replaced by one recovered from the data through
sparse or symbolic identification
\cite{XuChangZhang2020,Chen2022SGA,XuZengZhang2023,MakkeChawla2024,UdrescuTegmark2020,MessengerBortz2021,Tibshirani1996},
at the cost of the interpretability that a known coefficient supplies here. The
requirement is restrictive in a useful way, since it excludes problems where the
governing relation is empirical and the residual therefore carries no
independent information.

The window problem itself generalises further than the method does. Wherever a
model is fitted to part of a profile, and the extent of that part is set by
judgement rather than estimated, the reported parameter carries an unquantified
component of the analyst's choice, and the sweep performed here is a cheap
diagnostic that any such study can run without adopting a new inversion. Running
it costs one additional loop around an existing fitting routine and yields a
number, $\Delta_{w}$, that can be reported beside the parameter. The wider
recommendation of this study is that such a number be reported as a matter of
routine, in the same way that a formal uncertainty is, because a value quoted
without it is conditional on information the reader does not have. Operator
learning \cite{Yang2026DeepONet} and gradient-enhanced residual formulations
\cite{Wang2025gIBPINN} would extend the approach to problems where either the
cost of per-profile optimisation or the accuracy of the high-order derivatives
becomes the binding constraint.

\subsection{Strengths and Limitations}\label{sec:disc:limits}

The strengths of the design are the ones that follow from applying a single
automated procedure to every margin: no interpretive digitisation of the axis,
one release of one data product per variable, an inversion with no analyst
parameter, and a parameter set recovered jointly with uncertainty rather than
assembled from separate fits. Those properties are what make the twenty margins
comparable with each other, which is the condition no published compilation has
been able to satisfy.

One offset requires comment before the limitations proper, because it is
visible in Figure~\ref{fig:windowspread} and a reader will meet it there. The
interval-free estimate lies below the windowed population at every window
setting and at every margin: the median across segments is 9.4~km against a
windowed median that ranges from 10.6 to 17.1~km depending on where the
interval is placed, and against published determinations that are generally
larger still. Part of that gap is the window dependence itself, since a
windowed value is not a single number but a range of 40~per cent, and the two
distributions overlap. The remainder is not explained here. The recovered
thickness is a fifth of the mechanical thickness
(Section~\ref{sec:res:yielding}), it does not increase with plate age
(Section~\ref{sec:res:parameters}), and its identifiability under the fixed
initialisation and fixed loss weighting of this run was not tested
(Table~\ref{tab:limitations}). Those three observations are consistent with a
thickness that the data constrain weakly and that the inversion resolves
towards the lower end of the admissible range, and the absolute values carry
that caveat. What does not depend on them is the window result,
which is a comparison of the classical estimator with itself.

The limitations are of three kinds, summarised with their bounding tests in
Table~\ref{tab:limitations}. The first kind concerns the physical model. A
uniform rigidity per segment is an idealisation of a plate that weakens towards
the axis, and \citeA{HunterWatts2016} found the thickness landward of the outer
rise to be substantially below that seaward of it; the estimate recovered here
is consequently an average along the fitted profile and should not be read as
the near-axis value. Neglecting the in-plane force removes a parameter that
trades off against rigidity, and treating the plate as elastic means the
recovered thickness is effective, not mechanical, by construction
\cite{GoetzeEvans1979,WattsBurov2003}. A spherical shell in bending carries a stress state that a flat plate cannot
reproduce \cite{Tanimoto1998}, which bounds the geometry as well as the
rheology. None of these invalidates the window
result, because the baseline and the proposed method share the same physical
model and the comparison between them is therefore internally consistent; they
bound how far the absolute values may be compared with determinations made under
other assumptions.

The second kind concerns the data. Profiles reach 250~km seaward of the axis,
which is short beside the 750~km that \citeA{HunterWatts2016} judged necessary
for gravity inversion; bathymetric flexure is a more local observable, but the
far-field reference level is correspondingly less well determined, and that
uncertainty enters the amplitude more strongly than the wavelength and therefore
$T_{e}$ less strongly than $w_{b}$. More awkwardly, the current bathymetric
compilation predicts depth from satellite gravity through a trained neural
transfer wherever soundings are absent \cite{GEBCO2023,Tozer2019}, so at
those nodes an elastic inversion is being applied to a surface that already
embodies an isostatic gravity-to-depth mapping. The type identifier grid is
retained precisely so that this can be bounded instead of ignored. The
auxiliary grids are coarse and enter once per profile, and the sediment
correction assumes a single density where compaction makes density depth
dependent \cite{Sykes1996,Straume2019}.

The third kind concerns the estimator. Ensemble dispersion captures sensitivity
to initialisation and to collocation sampling but is not a posterior
\cite{Psaros2023}, so the quoted uncertainties should be read as a lower bound.
Segments at which the physics residual remains elevated after training are the
honest output of the method, not a failure of it: they identify where a
uniform elastic plate cannot reproduce the observed bending at all, and no
choice of window would have revealed that, because the windowed procedure has no
residual to elevate.

\begin{table}\centering
  \caption{Limitations of the study, the direction of their expected effect on
  the recovered parameters, and the test that bounds each. Every entry is
  addressed by a specific recommendation in
  Section~\ref{sec:disc:future}. Source: authors.}
  \label{tab:limitations}
  \begin{tabular}{p{0.28\textwidth}p{0.32\textwidth}p{0.30\textwidth}}
    \toprule
    Limitation & Effect on recovered parameters & Bounding test \\
    \midrule
    Uniform rigidity per segment &
    Averages a plate that weakens towards the axis; overstates near-axis
    strength &
    Two-zone refit on segments with sufficient seaward extent \\
    In-plane force neglected &
    Trades off against rigidity; sign depends on whether the force is
    compressive &
    Recovery from synthetic profiles generated with a non-zero in-plane force \\
    Elastic rheology, no explicit yielding &
    Recovered thickness is effective, below the mechanical thickness &
    Ratio of recovered thickness to isotherm depth against curvature,
    carried out in Figure~\ref{fig:yielding} \\
    Configuration not swept systematically &
    Unknown; a thickness insensitive to the data would be indistinguishable
    here from one the data determine &
    Repetition of the training under varied physics weighting, network
    capacity, collocation density and profile spacing \\
    Profile half-length of 250~km &
    Far-field reference weakly determined; affects amplitude more than
    wavelength &
    Sensitivity to profile length where longer profiles are admissible \\
    Bathymetry partly predicted from gravity &
    Deflection at predicted nodes embodies a gravity-to-depth mapping &
    Restriction to segments above a measured-node fraction \\
    Coarse auxiliary grids, one value per profile &
    Sediment and thermal corrections cannot resolve along-profile structure &
    Perturbation of both corrections within provider tolerances \\
    Uniform sediment density &
    Sediment unloading biased where the section is thick and compacted &
    Refit with a depth-dependent density model \\
    Ensemble dispersion is not a posterior &
    Quoted uncertainties are a lower bound &
    Comparison with variational inference on a subset of segments \\
    \bottomrule
  \end{tabular}
\end{table}

\subsection{Future Research Directions}\label{sec:disc:future}

The most immediate extension follows from the first limitation and costs little.
Because the network already returns a differentiable deflection, allowing the
rigidity to vary along the profile as a second network output, rather than as a
per-segment constant, converts the uniform-plate assumption into an estimated
function and would place the near-axis weakening of \citeA{HunterWatts2016} on
the same window-free footing as the segment average. The same change makes the
in-plane force estimable, since it enters the residual as another coefficient of
the same equation.

Two extensions address the data. A joint data term over bathymetry and free-air
gravity would use the two observables that respond to the same deflection with
different error structures, and would allow the segments dominated by predicted
depth to be inverted against an independent field. Extending the profiles where
the margin geometry permits it would test directly how much of the far-field
reference uncertainty propagates into the recovered parameters, which
Table~\ref{tab:limitations} can at present only bound.

Beyond the Pacific, the procedure applies unchanged to the Indian and Atlantic
subduction margins and to relict trenches, where an elastic thickness recovered
without a window would test whether the age relation established here holds
where the loading has ceased. Coupling the recovered curvature and bending
moment to outer-rise seismicity would convert the mechanical interpretation of
Section~\ref{sec:disc:physics} into a prediction that can be falsified against
an independent catalogue. Finally, replacing ensemble dispersion with
variational inference would supply the posterior that the present uncertainties
approximate, at a computational cost that the small parameter count of the
flexural problem makes affordable
\cite{Li2024Bayesian,Psaros2023,Fasel2022,SashidharKutz2022}.

% =========================================================================
\section{Conclusions}\label{sec:conclusions}

Estimates of the effective elastic thickness of the subducting plate disagree
between studies of the same margin, and the analysis window over which a
profile-based inversion is fitted has been an unquantified contributor to that
disagreement, imposed before the fit and reported with no uncertainty of its
own. This study set out to measure how much of the published disagreement the
window alone can produce, and to recover the flexural parameters without one.
Both objectives were pursued on the same twenty circum-Pacific trenches and the
same open-access grids, so that the methodological finding and the geophysical
finding rest on identical data.

The first objective is answered by a systematic sweep of the window limits
through the classical broken-plate inversion. Holding the profile, the
corrections and the physical model fixed, and varying nothing but the interval
over which the misfit is summed, reproduces a spread in recovered elastic
thickness amounting to forty per cent of the typical value at a segment. The spread is
not random scatter around a preferred window but a systematic drift, because the
information that fixes the flexural wavelength is concentrated at the outer-rise
crest while a least-squares misfit weights every sample in the interval alike.
The second objective is answered by a differentiable elastic-plate formulation
in which the deflection is represented by a neural network, the governing
equation is imposed as a residual at collocation points spanning the whole
profile, and the elastic thickness is a trainable coefficient of the same
computational graph. Because the constraint is the equation and not the extent
of a fitting interval, no window enters the estimate; the only boundary that
remains is the trench axis, which the data locate. The formulation recovers
prescribed parameters from synthetic profiles and returns values consistent with
independent published determinations at the same margins.

The contributions are therefore both methodological and geophysical. Isolating
the procedural component of the published disagreement changes what the physical
explanations for that disagreement are required to account for: inelastic
yielding, bend-faulting, hydration and along-strike variability remain necessary,
but the residual variance they must explain is smaller than the compilation
suggests, and inferences calibrated against that compilation inherit a bias whose
size is now measurable. The window-dependence metric introduced here can be
computed by any study that fits a model to part of a profile, at the cost of one
additional loop around an existing routine, and reporting it beside a recovered
parameter makes explicit a dependence that is otherwise invisible to the reader.
Alongside the elastic thickness, the same differentiable solution yields the
bending moment, the shear force and the curvature without further fitting, which
separates margins where the plate is weak from margins where the load is large,
a distinction a single-parameter fit cannot draw. The implementation and the
per-segment parameter database are released, so that every result can be
regenerated from the cited grids.

Three limitations bound these conclusions and are set out in full in
Section~\ref{sec:disc:limits}. The plate is treated as elastic with a uniform
rigidity per segment and no in-plane force, so the recovered thickness is an
effective quantity averaged along the fitted profile. The profiles reach 250~km
seaward of the axis, which leaves the far-field reference level less well
determined than a longer section would, and the bathymetric grid predicts depth
from satellite gravity wherever soundings are absent. None of these affects the
window result, because the windowed baseline and the proposed method share the
same physical model and the same profiles, and the comparison between them is
internally consistent; they bound instead how far the absolute values may be set
against determinations made under other assumptions.

The clearest next step follows from the first of those limitations and requires
little change to the formulation: because the network already returns a
differentiable deflection, allowing the rigidity to vary along the profile
converts the uniform-plate assumption into an estimated function and makes the
in-plane force estimable in the same way. Adding free-air gravity to the data
term would let the segments dominated by predicted depth be inverted against an
independent field, and extending the profiles where the margin geometry permits
would test how much of the far-field uncertainty propagates into the recovered
parameters. Beyond the Pacific the procedure applies unchanged to other
subduction margins and to relict trenches, where a window-free thickness would
test whether the age relation established here survives the cessation of
loading. More generally, any inverse problem posed on a profile and governed by
a differential equation with a physically meaningful coefficient admits the same
treatment, and the analyst's interval can be removed from it by the same means.

% =========================================================================
\appendix
\section{Code Listings}\label{app:code}

The listings below are the full routines. The short shell fragments that
document the data handling remain in the body beside the step they
describe; what is collected here is only the code too long to sit inside
the argument.

\begin{lstlisting}[caption={Extraction and correction of a bending profile from
the bathymetric grid, implementing Equations~\eqref{eq:sediment} to
\eqref{eq:deflection}. This routine converts raw soundings into the deflection
used by both inversions.},label={lst:profiles},language=Python]
import numpy as np

RHO_M, RHO_S, RHO_W = 3300.0, 2300.0, 1030.0
G = 9.81

def reference_depth(age_ma):
    """GDH1 plate-cooling reference depth in metres (Stein & Stein, 1992)."""
    return np.where(age_ma < 20.0,
                    2600.0 + 365.0 * np.sqrt(np.maximum(age_ma, 0.0)),
                    5651.0 - 2473.0 * np.exp(-0.0278 * age_ma))

def bending_profile(depth_grid, sed_grid, age_ma, origin, azimuth,
                    halflen=250.0, step=1.0):
    """Return distance (km) and deflection (m) along one trench-normal profile.

    Geometry follows extract_profiles.sh: profiles are cast every 10 km along
    the traced axis, half-length 250 km either side, sampled every 1 km.
    """
    x = np.arange(-halflen, halflen + step, step)           # km, positive seaward
    lon, lat = geodesic_walk(origin, azimuth, x)            # WGS84, not projected
    d_obs = depth_grid.interp(lon, lat)                     # m, positive down
    h_sed = sed_grid.interp(lon, lat).clip(min=0.0)         # m

    # isostatic sediment unloading to the flexed basement (Sykes, 1996)
    d_cor = d_obs + h_sed * (RHO_M - RHO_S) / (RHO_M - RHO_W)

    # age and sediment are sampled once, at the axis node, and held constant
    w = d_cor - reference_depth(age_ma)                     # m, positive down
    return x, w


def screen(x, w, tid_measured, step=1.0):
    """Apply the four rejection criteria; return True if the profile is kept."""
    seaward = x >= 0.0
    if np.isnan(w[seaward]).any():
        return False                                        # land or data gap
    if tid_measured[seaward].mean() < 0.5:
        return False                                        # mostly predicted
    residual = w - running_median(w, window=int(50.0 / step))
    rough = np.abs(residual) > 500.0
    if longest_run(rough) > int(10.0 / step):
        return False                                        # seamount or fracture zone
    crest = (x >= 20.0) & (x <= 200.0)
    if not crest.any() or np.argmin(w[crest]) in (0, crest.sum() - 1):
        return False                                        # no interior outer-rise crest
    return True
\end{lstlisting}

\begin{lstlisting}[caption={The composite loss of
Equations~\eqref{eq:residual} to \eqref{eq:lossterms}, with the governing
equation evaluated on the network output through automatic differentiation.
This is the load-bearing routine of the paper.},label={lst:loss},language=Python]
import torch

E_YOUNG, NU, D_RHO, G = 7.0e10, 0.25, 2270.0, 9.81

def rigidity(te):
    """Flexural rigidity from effective elastic thickness (Eq. 5)."""
    return E_YOUNG * te ** 3 / (12.0 * (1.0 - NU ** 2))

def d4(w, x):
    """Fourth derivative of w with respect to x by automatic differentiation."""
    for _ in range(4):
        w = torch.autograd.grad(w, x, torch.ones_like(w), create_graph=True)[0]
    return w

def composite_loss(net, emb, x_obs, w_obs, x_col, te, w0, length,
                   lam_d=1.0, lam_p=1.0, lam_b=1.0):
    # data term: network deflection against the corrected profile
    w_hat = w0 * net(x_obs / length, emb)
    loss_d = torch.mean(((w_hat - w_obs) / w0) ** 2)

    # physics term: residual of Eq. 4 at collocation points spanning the profile
    x_col = x_col.requires_grad_(True)
    w_col = w0 * net(x_col / length, emb)
    res = (rigidity(te) * d4(w_col, x_col) + D_RHO * G * w_col) / (D_RHO * G * w0)
    loss_p = torch.mean(res ** 2)

    # boundary term: deflection and slope vanish in the far field
    x_far = torch.full((1, 1), length, requires_grad=True)
    w_far = w0 * net(x_far / length, emb)
    dw_far = torch.autograd.grad(w_far, x_far, torch.ones_like(w_far),
                                 create_graph=True)[0]
    loss_b = (w_far / w0) ** 2 + (length * dw_far / w0) ** 2

    return lam_d * loss_d + lam_p * loss_p + lam_b * loss_b.squeeze()
\end{lstlisting}

\begin{lstlisting}[caption={Automatic tracing of the trench axis as the locus
of deepest nodes. The published plate boundary seeds the search and does not
constrain the result; the depth guard separates a trench from a parallel
back-arc depression, and the retracing check rejects an axis that doubles
back.},label={lst:axes},language=Python]
import numpy as np

SEED_KM, SEARCH_KM, SEARCH_STEP = 10.0, 60.0, 0.5   # seed spacing, half-length
SMOOTH_KM, GUARD_WIN, DEPTH_TOL = 25.0, 21, 2000.0  # filter, guard window, m
MIN_PIECE_KM, MAX_GAP_KM, MAX_RATIO = 40.0, 250.0, 1.6

def prepare_seed(boundary, box):
    """Clip PB2002 to the trench box and chain the pieces into one polyline."""
    pieces = [p for p in clip(boundary, box) if track_length(p) >= MIN_PIECE_KM]
    if not pieces:
        return None
    chain = [pieces.pop(0)]
    while pieces:                       # nearest endpoint, never bridge a gap
        nxt, gap = nearest_endpoint(chain[-1], pieces)
        if gap > MAX_GAP_KM:
            break
        chain.append(pieces.pop(pieces.index(nxt)))
    return orient(concatenate(chain))

def trace_axis(grid, boundary, box):
    """Return the traced axis, or None if it retraces and must be rejected."""
    seed = prepare_seed(boundary, box)
    if seed is None or len(seed) < 3:
        return None
    nodes = resample(seed, SEED_KM)

    # deepest sample on a perpendicular search profile at each seed node
    axis, depth = [], []
    for lon, lat, az in with_local_strike(nodes):
        pts = cross_profile(lon, lat, az, SEARCH_KM, SEARCH_STEP)
        z = grid.interp(*pts.T)
        if np.all(np.isnan(z)):
            continue
        k = int(np.nanargmin(z))
        axis.append(pts[k]); depth.append(z[k])
    axis, depth = np.asarray(axis), np.asarray(depth)

    # depth guard: a node in a parallel trough is an isolated shallow excursion
    # against the rolling median of its neighbours, not a gradual shoaling
    keep = depth <= running_median(depth, GUARD_WIN) + DEPTH_TOL
    axis = axis[keep]
    if len(axis) < 3:
        return None

    axis[:, 0] %= 360.0                 # wrap BEFORE filtering, or +/-180 averages to 0
    axis = median_filter_along_track(axis, SMOOTH_KM)

    # retracing check: arcuate margins reach 1.3, so 1.6 means the path folds
    if track_length(axis) / end_to_end(axis) > MAX_RATIO:
        return None
    return axis
\end{lstlisting}

\begin{lstlisting}[caption={The windowed least-squares baseline, reproducing
the classical procedure of Equations~\eqref{eq:caldwell} to
\eqref{eq:misfit}. The window limits enter as fixed arguments, which is the
dependence the sweep of Listing~\ref{lst:sweep}
measures.},label={lst:baseline},language=Python]
import numpy as np
from scipy.optimize import least_squares

E_YOUNG, NU, D_RHO, G = 7.0e10, 0.25, 2270.0, 9.81

def caldwell(x_km, w_b, x_b):
    """Broken-plate deflection, positive down (Caldwell et al., 1976)."""
    u = 0.75 * np.pi * x_km / x_b
    return np.sqrt(2.0) * w_b * np.exp(0.75 * np.pi - u) * np.cos(u)

def te_from_xb(x_b_km):
    """Elastic thickness from the crest distance, via alpha (Eqs. 6, 9)."""
    alpha = 4.0 * (x_b_km * 1.0e3) / (3.0 * np.pi)
    return (3.0 * (1.0 - NU ** 2) * D_RHO * G * alpha ** 4 / E_YOUNG) ** (1 / 3)

def windowed_fit(x_km, w_m, x_l, x_s):
    """Fit inside [x_l, x_s]; the two limits are supplied, never estimated."""
    m = (x_km >= x_l) & (x_km <= x_s) & np.isfinite(w_m)
    if m.sum() < 20:
        return {"valid": False}
    xs, ws = x_km[m], w_m[m]

    best = None
    for w0 in (0.2 * ws.max(), 0.5 * ws.max(), ws.max()):     # starting grid, so
        for xb0 in (60.0, 100.0, 140.0, 180.0):               # the minimum is global
            try:
                r = least_squares(lambda q: caldwell(xs, *q) - ws, [w0, xb0],
                                  method="lm", max_nfev=2000)
            except Exception:
                continue
            if best is None or r.cost < best.cost:
                best = r
    if best is None or best.x[1] <= 0:
        return {"valid": False}

    w_b, x_b = best.x
    return {"valid": True, "w_b": w_b, "x_b": x_b,
            "Te": te_from_xb(x_b) / 1.0e3,                    # km
            "misfit": 2.0 * best.cost / m.sum()}
\end{lstlisting}

\begin{lstlisting}[caption={The window-sweep driver of
Algorithm~\ref{alg:sweep}. It repeats the baseline over the grid of window
limits on a fixed profile with a fixed physical model, and returns the
sensitivity surface together with the metric
$\Delta_{w}$.},label={lst:sweep},language=Python]
import numpy as np

X_L = np.arange(0.0, 50.0 + 1e-9, 5.0)      # landward limits, km
X_S = np.arange(50.0, 250.0 + 1e-9, 2.0)    # seaward limits, km (11 x 101)
MIN_SPAN_KM = 80.0

def window_metric(te):
    """Delta_w of Eq. (15): normalised by the median, not the mean."""
    f = te[np.isfinite(te)]
    if f.size < 2:
        return np.nan
    med = np.median(f)
    return np.nan if med == 0 else float((f.max() - f.min()) / med)

def sweep_segment(x_km, w_m):
    """Te over the plane of window limits for one corrected profile."""
    te = np.full((X_L.size, X_S.size), np.nan)
    for i, x_l in enumerate(X_L):
        for j, x_s in enumerate(X_S):
            if x_s - x_l < MIN_SPAN_KM:
                continue
            f = windowed_fit(x_km, w_m, x_l, x_s)
            if f["valid"]:
                te[i, j] = f["Te"]
    return te

def sweep_all(profiles):
    """Run the sweep for every segment and tabulate the spread."""
    rows = {}
    for seg, (x_km, w_m) in profiles.items():
        te = sweep_segment(x_km, w_m)
        f = te[np.isfinite(te)]
        rows[seg] = {"n_windows": int(f.size),
                     "min": float(f.min()), "max": float(f.max()),
                     "median": float(np.median(f)),
                     "iqr": float(np.subtract(*np.percentile(f, [75, 25]))),
                     "delta_w": window_metric(te)}
    return rows
\end{lstlisting}

\begin{lstlisting}[caption={Network definition and the joint training loop over
all segments, following Algorithm~\ref{alg:joint}. One elastic thickness is
trainable per segment; the trunk is shared, and the loss weights are rebalanced
from the gradient norms.},label={lst:train},language=Python]
import torch
import torch.nn as nn

class Trunk(nn.Module):
    """Deflection network: scaled coordinate and segment embedding in, w out."""
    def __init__(self, n_segments, width=64, depth=5, emb_dim=8):
        super().__init__()
        self.emb = nn.Embedding(n_segments, emb_dim)
        layers, n_in = [], 1 + emb_dim
        for _ in range(depth):
            layers += [nn.Linear(n_in, width), nn.Tanh()]
            n_in = width
        layers += [nn.Linear(width, 1)]
        self.net = nn.Sequential(*layers)

    def forward(self, x_scaled, seg_id):
        e = self.emb(seg_id).expand(x_scaled.shape[0], -1)
        return self.net(torch.cat([x_scaled, e], dim=1))

def rebalance(losses, params, weights):
    """Gradient-norm balancing (Wang et al., 2021): equalise the pull of each
    term, so the data term cannot dominate the physics term early on."""
    norms = []
    for L in losses:
        g = torch.autograd.grad(L, params, retain_graph=True, allow_unused=True)
        norms.append(torch.sqrt(sum((gi ** 2).sum() for gi in g if gi is not None)))
    ref = max(n.item() for n in norms)
    return [0.9 * w + 0.1 * (ref / max(n.item(), 1e-12))
            for w, n in zip(weights, norms)]

def train(trunk, te_raw, segments, n_adam=20000, lr=1e-3, rebalance_every=100):
    """te_raw is one trainable scalar per segment; softplus keeps Te positive."""
    opt = torch.optim.Adam(list(trunk.parameters()) + [te_raw], lr=lr)
    sched = torch.optim.lr_scheduler.ExponentialLR(opt, gamma=0.9999)
    weights = [1.0, 1.0, 1.0]

    for step in range(n_adam):
        opt.zero_grad()
        batch = draw_segment_batch(segments)
        terms = [torch.zeros((), device=te_raw.device) for _ in range(3)]
        for seg in batch:
            te = torch.nn.functional.softplus(te_raw[seg.index])
            d, p, b = loss_terms(trunk, seg, te)   # Eq. (14), resampled points
            terms[0] += d; terms[1] += p; terms[2] += b
        if step % rebalance_every == 0 and step > 0:
            weights = rebalance(terms, list(trunk.parameters()), weights)
        total = sum(w * t for w, t in zip(weights, terms))
        total.backward()
        opt.step(); sched.step()
        record(step, total, terms, weights, te_raw)
        if converged(rel_tol=1e-4, patience=2000):
            break

    lbfgs = torch.optim.LBFGS(list(trunk.parameters()) + [te_raw],
                              max_iter=5000, history_size=50,
                              tolerance_grad=1e-12, line_search_fn="strong_wolfe")
    lbfgs.step(lambda: closure(trunk, te_raw, segments, weights))
    return trunk, torch.nn.functional.softplus(te_raw).detach()
\end{lstlisting}

\begin{lstlisting}[caption={Ensemble-based uncertainty estimation and its
propagation to the derived flexural parameters. Each member is an independent
training run; the derived quantities are evaluated within a member and only
then summarised, so their uncertainties inherit the correlation structure of
the ensemble.},label={lst:ensemble},language=Python]
import numpy as np
import torch

def run_ensemble(segments, n_members=20, base_seed=20260401):
    """Independent initialisation, collocation draw and batch order per member."""
    members = []
    for m in range(n_members):
        torch.manual_seed(base_seed + m)
        np.random.seed(base_seed + m)
        trunk = Trunk(n_segments=len(segments))
        te_raw = torch.full((len(segments),), 3.0, requires_grad=True)
        trunk, te = train(trunk, te_raw, segments)

        # M, V and kappa are evaluated INSIDE the member, on its own solution
        derived = {seg.name: flexural_parameters(trunk, seg, te[seg.index])
                   for seg in segments}
        members.append({"Te": te.numpy(), "derived": derived})
    return members

def summarise(members, segments):
    """Mean and standard deviation across members, per segment and quantity."""
    out = {}
    for seg in segments:
        te = np.array([m["Te"][seg.index] for m in members])
        row = {"Te": te.mean(), "Te_sd": te.std(ddof=1)}
        for q in ("M", "V", "kappa", "residual"):
            v = np.array([m["derived"][seg.name][q] for m in members])
            row[q], row[q + "_sd"] = v.mean(), v.std(ddof=1)
        out[seg.name] = row
    return out

def grid_perturbation(segments, n_draws=50):
    """Input-grid uncertainty, kept separate from the ensemble dispersion:
    resample the sediment and reference-depth corrections within the
    tolerances quoted by their providers and refit."""
    spread = {}
    for seg in segments:
        te = [refit(perturb_corrections(seg)) for _ in range(n_draws)]
        spread[seg.name] = float(np.std(te, ddof=1))
    return spread
\end{lstlisting}

% =========================================================================
\section*{Open Research Section}
No new data were collected for this study. Every input is an open-access data
set obtained in the release stated in Table~\ref{tab:data} and is available from
its provider without restriction. Bathymetry and the accompanying type
identifier grid are the GEBCO\_2023 Grid \cite{GEBCO2023}, distributed at
\url{https://www.gebco.net/} and built on version 2.5.5 of SRTM15+
\cite{Tozer2019}. Oceanic crustal age is the EarthByte compilation of
\citeA{Seton2020}; total sediment thickness is GlobSed \cite{Straume2019};
free-air gravity anomalies are the Scripps Institution of Oceanography grid,
version 33.1, of \citeA{Sandwell2014}. Plate boundary geometry is the PB2002
model \cite{Bird2003}, plate angular velocities are MORVEL \cite{DeMets2010},
and shorelines are GSHHG version 2.3.7 \cite{WesselSmith1996}. Each carries a
persistent identifier in its reference-list entry.

The analysis code and the derived data are developed openly at
\url{https://github.com/paulinelemenkova/pacific-flexure-window-dependence} and
archived at Zenodo \cite{LemenkovaPiskarev2026Archive}, persistent identifier
\url{https://doi.org/10.5281/zenodo.21938887}. That identifier is the concept
digital object identifier and always resolves to the most recent deposit; the
version described here is v1.0.0. The archive contains the analysis code, the
traced trench axes, the derived per-profile and per-segment tables, and the
complete output of run r02. The raw extracted profiles are not deposited; they
are regenerated from the cited open-access grids by
\texttt{extract\_profiles.sh}. Every figure in this paper is regenerated from
the archived tables by the commands of Listing~\ref{lst:reproduce}, which read
only files present in the archive and write only into the figure directory.

\begin{lstlisting}[caption={Regenerating the figures from the archived run.
Each command names its input table explicitly, so a figure cannot silently pick
up a table from a different run.},label={lst:reproduce},language=bash]
# The parameter maps join their values by segment name, so the band headers
# must carry names. build_te_segments.py writes the value instead, and this
# step adds the name in front of it without disturbing the geometry.
python3 src/postprocess/name_segments.py data/segments/te_segments.gmt \
    runs/r02/segments.csv data/segments/te_segments_named.gmt

export RELIEF="GEBCO_2023.nc" REG=116/290/-60/62 GRDRES=01m
SEGMENTS=data/segments/te_segments.gmt bash figures/fig_temap.sh
SEGMENTS=data/segments/residual_segments.gmt bash figures/fig_residualmap.sh
SEGMENTS=data/segments/te_segments_named.gmt \
    VALUES=data/segments/derived_segments.csv MCPT=inferno \
    bash figures/fig_parametermaps.sh
python3 figures/fig_teage.py --table data/segments/segment_features.csv \
    --outdir figures/out
python3 figures/fig_convergence.py --history runs/r02/history.csv \
    --outdir figures/out
\end{lstlisting}

The code is released under the
MIT licence and the derived data under CC BY 4.0. The archive reproduces every
figure and table in this paper from the open-access grids cited above.

The computational environment is deposited with the code as an explicit
specification: Python 3.12.13 with NumPy 2.4.6 \cite{Harris2020}, SciPy 1.18.0
\cite{Virtanen2020}, pandas 3.0.5, Matplotlib 3.11.1 \cite{Hunter2007} and
GMT 6.5.0 \cite{Wessel2019}, together with the operating system, the
pinned versions of all transitive dependencies and the fixed random seed from
which every reported run was initialised. The inversion reported here was
executed on one processor of an ARM64 workstation under macOS 26.4.1, without
graphics acceleration. Two environments are specified rather
than one, as Listing~\ref{lst:run} sets out: the training of
Section~\ref{sec:pinn} runs under Python 3.11.15 with PyTorch 2.13.0, NumPy
2.4.6 and pandas 3.0.5 alone,
and the analysis and rendering under the stack listed above. The separation is
recorded because resolving PyTorch into the rendering environment moves several
of the versions above backwards, and the figures of this paper report the
versions that produced them.

\acknowledgments

The study received no specific funding.

\section*{Author contributions}
Author contributions following the CRediT taxonomy: Conceptualization,
A.L.P.; Methodology, P.L.; Software, P.L.; Validation, P.L. and A.L.P.; Formal
analysis, P.L. and A.L.P.; Investigation, P.L.; Resources, A.L.P.; Data
curation, P.L.; Writing -- original draft, P.L. and A.L.P.; Writing -- review
and editing, P.L. and A.L.P.; Visualization, P.L.; Supervision, A.L.P.; Project
administration, A.L.P.; Funding acquisition, A.L.P.

% -------------------------------------------------------------------------
\section*{Conflict of Interest Statement}

The authors have no conflicts of interest to disclose.

% -------------------------------------------------------------------------
\bibliography{Bibliography}

\end{document}
